\documentclass[12pt, letterpaper, oneside]{article}

% --- PAQUETES DE IDIOMA Y FORMATO ---
\usepackage[spanish, es-tabla]{babel}
\usepackage[margin=2.5cm]{geometry}
\usepackage{setspace}
\onehalfspacing % Interlineado 1.5 para formato universitario en 1 columna
\usepackage[utf8]{inputenc} % Permite escribir acentos y caracteres UTF-8 directamente
\usepackage[T1]{fontenc}    % Asegura la correcta renderización y salida de caracteres con acentos
\usepackage[spanish]{babel} % Configura la tipografía e hifenación en español

% --- PAQUETES PARA MATEMÁTICAS AVANZADAS Y CRIPTOGRAFÍA ---
\usepackage{amsmath, amssymb, amsfonts}
\usepackage{float}

% --- PAQUETES PARA FIGURAS Y DIAGRAMAS ---
\usepackage{graphicx}
\graphicspath{{figuras/}}
\usepackage{tikz}
\usetikzlibrary{shapes.geometric, arrows}

% =====================================================
% Paquetes para el capítulo de benchmarking criptográfico
% =====================================================
\usepackage{booktabs}
\usepackage{tabularx}
\usepackage{longtable}
\usepackage{multirow}
\usepackage{array}
\usepackage{siunitx}
\usepackage{graphicx}
\usepackage{float}
\usepackage{xcolor}
\usepackage{pgfplots}
\usepackage{tikz}
\usepackage{amsmath}
\usepackage{amssymb}

\pgfplotsset{compat=1.18}

% Columnas útiles para tablas
\newcolumntype{Y}{>{\raggedright\arraybackslash}X}
\newcolumntype{C}{>{\centering\arraybackslash}X}
\newcolumntype{R}{>{\raggedleft\arraybackslash}X}

% Marcador visual para información pendiente
\newcommand{\todoresult}[1]{\textbf{[TODO: #1]}}

% --- PAQUETE PARA BLOQUES DE CÓDIGO FUENTE ---
\usepackage{listings}
\usepackage{xcolor}

\definecolor{codegray}{rgb}{0.5,0.5,0.5}
\definecolor{codepurple}{rgb}{0.58,0,0.82}
\definecolor{backcolour}{rgb}{0.95,0.95,0.92}

\lstdefinestyle{estiloCodigo}{
    backgroundcolor=\color{backcolour},   
    commentstyle=\color{codegray},
    keywordstyle=\color{blue},
    numberstyle=\tiny\color{codegray},
    stringstyle=\color{codepurple},
    basicstyle=\ttfamily\footnotesize,
    breakatwhitespace=false,         
    breaklines=true,                 
    captionpos=b,                    
    keepspaces=true,                 
    numbers=left,                    
    numbersep=5pt,                  
    showspaces=false,                
    showstringspaces=false,
    showtabs=false,                  
    tabsize=2
}
\lstset{style=estiloCodigo}

% --- CONFIGURACIÓN DE NÚMEROS Y ENLACES ---
\usepackage{hyperref}
\hypersetup{
    colorlinks=true,
    linkcolor=black,
    citecolor=blue,
    urlcolor=blue
}

% --- INICIO DEL DOCUMENTO ---
\begin{document}

% 1. Portada Institucional
%\begin{titlepage}
    \centering
    \onehalfspacing
    \includegraphics[width=0.45\textwidth]{logo_uniandes.png} 

    % Portada 1
    % Título del proyecto
    {\bfseries \Large Autenticación con Protocolo Zero-Knowledge-Proof basado en Learning with Errors \par}
    \vfill

    % Autores
    {\large \textbf{Autores:}}\\[0.2cm]
    {\large Cesar Alejandro Avellaneda}\\[0.1cm]
    {\large Juan Esteban Álvarez}
    \vfill
    % Afiliación e Institución
    {\bfseries Universidad de los Andes}\\
    Facultad de Ingeniería\\
    Departamento de Ingeniería de Sistemas y Computación\\
    Bogotá, Colombia\\[0.3cm]
    2026 % Modifica la fecha aquí cuando corresponda

\newpage
\newpage
    %Portada 2

    %Titulo
    {\bfseries \Large Autenticación con Protocolo Zero-Knowledge-Proof basado en Learning with Errors \par}
    \vfill
    %Autores
    {\small Tesis presentada por:}
    {\large Cesar Alejandro Avellaneda}\\[0.1cm]
    {\large Juan Esteban Álvarez}
    % Requisito
    \textit{Como requisito parcial para optar al título de:}\\
    \textbf{Ingeniero de Sistemas y Computación}
    \vfill
    % Tutores
    {\large \textbf{Tutores:}}\\[0.2cm]
    {\large Sandra Julieta Rueda Rodriguez}\\[0.1cm]
    {\large Milton Emigdio Quiroga Becerra}
    \vfill
    % Afiliación e Institución
    {\bfseries Universidad de los Andes}\\
    Facultad de Ingeniería\\
    Departamento de Ingeniería de Sistemas y Computación\\
    Bogotá, Colombia\\[0.3cm]
    2026 % Modifica la fecha aquí cuando corresponda
\end{titlepage}
\newpage

% 2. Numeración arábiga para el contenido
\pagenumbering{arabic}

% Indices y tablas de contenidos
% --- ÍNDICES Y TABLAS DE CONTENIDO ---
\cleardoublepage
\tableofcontents % Genera la Tabla de Contenidos automática

\cleardoublepage
\listoffigures % (Opcional) Índice automático de figuras/diagramas

\cleardoublepage
\listoftables % (Opcional) Índice automático de tablas

\cleardoublepage
% --- FIN DE ÍNDICES ---

% 4. Inclusión de las Secciones / Capítulos (Módulos)
\section{Introducción}

\subsection{Generalidades}
La autenticación de usuarios constituye uno de los problemas más complejos dentro de la seguridad de la información contemporánea. De acuerdo con múltiples análisis e informes del sector, entre el 67\% y el 80\% de los problemas globales de seguridad están directamente vinculados con la autenticación, la suplantación de credenciales y el robo de identidades \cite{identity_mfa}.

A pesar de la evolución en las tecnologías de verificación mediante contraseñas de un solo uso (\textit{One-Time Passwords}, OTP), esquemas de doble factor basados en los estándares de la alianza FIDO y técnicas biométricas, estas alternativas aún presentan limitaciones considerables \cite{fido_alliance}:
\begin{itemize}
    \item \textbf{Contraseñas tradicionales:} Constituyen un mecanismo poco adecuado para las capacidades cognitivas humanas y obligan al usuario a ceder la totalidad de su secreto al sistema verificador.
    \item \textbf{Autenticación biométrica:} Involucra información de identificación personal no modificable, lo que genera objeciones significativas en materia de privacidad y cumplimiento normativo .
    \item \textbf{Esquemas de doble factor en servicios financieros:} En escenarios como redes de cajeros automáticos (ATM), la combinación de la tarjeta y el PIN implica entregar la totalidad del secreto a terceros y operadores de redes intermedias (e.g., Cirrus, Visa, MasterCard), comprometiendo la soberanía de los datos del usuario.
\end{itemize}

\subsection{Contexto de la Situación}
El panorama actual de la seguridad de la información se encuentra condicionado por dos dinámicas críticas: la necesidad de proteger la identidad sin entregar credenciales sensibles y la inminente llegada de la computación cuántica 

Por una parte, la autenticación convencional exige transmitir secretos legibles o equivalentes a través de nodos intermedios. Las Pruebas de Conocimiento Cero (\textit{Zero-Knowledge Proofs}, ZKP) ofrecen una alternativa criptográfica que permite a una parte probadora ($P$) demostrar a un verificador ($V$) que posee cierta información privada sin revelar un solo bit de dicho secreto.

Por otra parte, la ejecución de algoritmos como el de Shor y Grover en computadores cuánticos amenaza la seguridad de los esquemas criptográficos tradicionales . Al transformar las funciones clásicas de una sola vía (como el logaritmo discreto, la factorización de enteros o el isomorfismo de grafos) en funciones fácilmente invertibles, la computación cuántica dejará inútiles los protocolos clásicos, haciendo imperativa la migración hacia primitivas resistentes a capacidades cuánticas de ataque (\textit{quantum-safe}).

\subsection{Justificación}
Frente a la vulnerabilidad de las credenciales tradicionales y la amenaza de adversarios cuánticos, surge la necesidad de contar con esquemas de autenticación que garanticen la soberanía de los datos del usuario y la robustez técnica frente a ataques clásicos y cuánticos.

Los problemas matemáticos basados en retículas (\textit{lattices}), en particular, el problema \textit{Learning with Errors} (LWE) con distribución Gaussiana discreta permite construir funciones de una sola vía respaldadas por una reducción de seguridad del peor caso al caso promedio (\textit{worst-case to average-case reduction}). Esta propiedad garantiza que resolver una instancia aleatoria del sistema sea tan difícil como los problemas de retículos más complejos en el peor caso, manteniendo el cálculo inverso computacionalmente inviable tanto para la arquitectura tradicional Von-Neumann como para ordenadores cuánticos.


La integración de las Pruebas de Conocimiento Cero (ZKP) con la solidez matemática de LWE atiende simultáneamente ambas problemáticas:
\begin{enumerate}
    \item Resuelve la exposición de secretos al verificar la identidad del usuario sin revelar sus credenciales.
    \item Garantiza la seguridad del canal y la autenticación a largo plazo contra supercomputadoras cuánticas gracias a la complejidad multidimensional del problema LWE.
\end{enumerate}

\subsection{Objetivos}

\subsubsection{Objetivo General}
Definir el funcionamiento, robustez, alcance y riesgos de un protocolo de autenticación de conocimiento cero que sea resistente ante ataques cuánticos, en contraste con protocolos tradicionales, con el propósito de mejorar, para el presente y futuro previsible, la seguridad en sistemas que tengan un fuerte componente de autenticación.

\subsubsection{Objetivos Específicos}
\begin{enumerate}
    \item Analizar los fundamentos matemáticos y las limitaciones de los protocolos \textit{Zero-Knowledge-Proof}.
    \item Analizar los fundamentos matemáticos y las limitaciones del problema \textit{Learning with Errors}.
    \item Diseñar e implementar un prototipo de protocolo que integre ambos conceptos con fines de autenticación segura.
    \item Explorar las particularidades necesarias para el correcto funcionamiento del protocolo planteado.
    \item Comparar los protocolos de autenticación convencionales más relevantes con el protocolo planteado.
\end{enumerate}

\section{Marco Teórico}

\subsection{Zero-Knowledge Proofs (ZKP)}

Los protocolos de prueba de conocimiento cero (\textit{Zero-Knowledge Proofs} o ZKPs) son una técnica criptográfica que permite verificar la veracidad de una información sin necesidad de revelar los datos en sí \cite{mulder2023trends}, en el caso particular de esta tesis para verificar a un usuario ante una entidad bancaria. En este proceso interactúan dos partes: un \textbf{probador} (\textit{prover}), que genera y demuestra la validez de una afirmación, y un \textbf{verificador} (\textit{verifier}), quien comprueba dicha afirmación sin obtener información adicional ni conocer los datos subyacentes.

Las ZKP pueden clasificarse principalmente en dos categorías \cite{mulder2023trends}:
\begin{itemize}
    \item \textbf{Interactivas:} donde el verificador envía múltiples retos aleatorios al probador hasta alcanzar un alto grado de certeza.
    \item \textbf{No interactivas:} donde la validación se realiza en una sola interacción (\textit{single shot}).
\end{itemize}

La diferencia fundamental radica en que las pruebas no interactivas sustituyen los retos aleatorios del verificador por un \textbf{valor de referencia común} (\textit{common reference value}), lo que permite que la prueba generada sea transferible y validada de forma independiente por terceros \cite{mulder2023trends}.

La utilidad principal de este paradigma consiste en que en un protocolo convencional de autenticación se tiene el verificador y el usuario, el verificador no confía en el usuario antes de que este se identifique y el usuario se asume que tiene confianza de que la entidad verificadora no va a filtrar su información personal y por ende puede entregarla para llevar a cabo la autenticación. En cambio, usando ZKP se elimina la posibilidad de filtración y la necesidad de confianza en la entidad.

\subsubsection{Fundamentos Matemáticos y Propiedades de las Pruebas de Conocimiento Cero}

\paragraph{Formulación Formal de Relaciones NP y Sistemas de Prueba Interactivos}
En la teoría de la complejidad computacional y la criptografía moderna, un sistema de prueba de conocimiento cero se formaliza mediante el concepto de lenguaje en NP definido sobre una relación binaria evaluable en tiempo polinomial \cite{goldwasser1989knowledge}. Sea $R \subseteq \{0,1\}^* \times \{0,1\}^*$ una relación NP. Se define el lenguaje asociado $L_R$ como el conjunto de instancias o afirmaciones para las cuales existe un testigo (\textit{witness}) que satisface la relación:
\begin{equation}
L_R = \{ x \in \{0,1\}^* \mid \exists \, w \in \{0,1\}^* \text{ tal que } (x, w) \in R \}
\end{equation}
Donde $x$ representa la declaración pública (por ejemplo, una clave pública o el valor de un hash) y $w$ representa el testigo privado o secreto (por ejemplo, una clave privada o una preimagen).

Un sistema de prueba interactivo para el lenguaje $L_R$ consta de un par de máquinas de Turing probabilísticas que interactúan mediante el intercambio de mensajes en tiempo polinomial (PPT): el \textbf{Proponente} ($\mathcal{P}$) y el \textbf{Verificador} ($\mathcal{V}$). El proponente $\mathcal{P}$ posee la instancia pública $x$ y el testigo secreto $w$, mientras que el verificador $\mathcal{V}$ solo conoce $x$. Al finalizar el protocolo de interacción, denotado por el flujo de transcripción $\langle \mathcal{P}(x, w), \mathcal{V}(x) \rangle$, el verificador emite una decisión binaria:
\begin{equation}
\langle \mathcal{P}(x, w), \mathcal{V}(x) \rangle \in \{1 \text{ (Aceptar)}, 0 \text{ (Rechazar)}\}
\end{equation}

\paragraph{Definición Matemática Rigurosa de las Propiedades Fundamentales}
Para que la interacción $(\mathcal{P}, \mathcal{V})$ constituya un sistema de prueba de conocimiento cero (\textit{Zero-Knowledge Proof} - ZKP), debe satisfacer de forma estricta tres propiedades fundamentales cuantificadas bajo un parámetro de seguridad $\lambda \in \mathbb{N}$ \cite{goldwasser1989knowledge, badiezadegan2026complexity}:

\begin{enumerate}
    \item \textbf{Completitud (\textit{Completeness}):} Garantiza que si la afirmación es verdadera y ambas partes siguen el protocolo honestamente, el verificador siempre aceptará la prueba. Formalmente, para todo par $(x, w) \in R$:
    \begin{equation}
    \Pr\left[ \langle \mathcal{P}(x, w), \mathcal{V}(x) \rangle = 1 \right] \ge 1 - \text{negl}(\lambda)
    \end{equation}
    Donde $\text{negl}(\lambda)$ denota una función infinitesimal (\textit{negligible}), tal que para todo polinomio $p(\cdot)$ existe un $\lambda_0$ donde $\text{negl}(\lambda) < 1/p(\lambda)$ para todo $\lambda > \lambda_0$. Si la probabilidad es exactamente igual a $1$, el esquema posee \textit{completitud perfecta}.

    \item \textbf{Solidez (\textit{Soundness}):} Garantiza que si la afirmación es falsa ($x \notin L_R$), ningún proponente deshonesto $\mathcal{P}^*$ (incluso con recursos computacionales ilimitados) puede convencer al verificador con una probabilidad significativa. Formalmente, para todo $x \notin L_R$ y cualquier proponente deshonesto $\mathcal{P}^*$:
    \begin{equation}
    \Pr\left[ \langle \mathcal{P}^*(x), \mathcal{V}(x) \rangle = 1 \right] \le \epsilon(\lambda) \le \text{negl}(\lambda)
    \end{equation}
    La cantidad $\epsilon(\lambda)$ se denomina el \textit{error de solidez} (\textit{soundness error}). En sistemas donde la solidez se garantiza únicamente contra proponentes acotados a tiempo polinomial (PPT), el sistema se denomina un \textit{argumento de conocimiento} (\textit{Argument of Knowledge}).

    \item \textbf{Conocimiento Cero (\textit{Zero-Knowledge}):} Garantiza que el verificador no aprende ninguna información adicional sobre el testigo $w$ más allá de la veracidad de la afirmación $x \in L_R$. Se formaliza mediante la existencia de un \textbf{Simulador} probabilístico en tiempo polinomial ($\mathcal{S}$) que, teniendo acceso únicamente a la instancia pública $x$ (y a la máquina del verificador, pero sin acceso a $w$), es capaz de generar una transcripción sintética que resulta indistinguible de la interacción real.
    
    Para cualquier verificador PPT deshonesto $\mathcal{V}^*$ y para todo $(x, w) \in R$, la vista del verificador $\text{View}_{\mathcal{V}^*}\left[ \langle \mathcal{P}(x, w), \mathcal{V}^*(x) \rangle \right]$ satisface:
    \begin{equation}
    \mathcal{S}^{\mathcal{V}^*}(x) \approx \text{View}_{\mathcal{V}^*}\left[ \langle \mathcal{P}(x, w), \mathcal{V}^*(x) \rangle \right]
    \end{equation}
    Según la noción de indistinguibilidad ($\approx$), el esquema se clasifica en:
    \begin{itemize}
        \item \textbf{Conocimiento Cero Perfecto ($\approx_p$):} Las distribuciones de probabilidad son matemáticamente idénticas.
        \item \textbf{Conocimiento Cero Estadístico ($\approx_s$):} La distancia estadística entre ambas distribuciones es infinitesimal.
        \item \textbf{Conocimiento Cero Computacional ($\approx_c$):} Ningún algoritmo PPT puede distinguir entre ambas distribuciones con una ventaja no despreciable.
    \end{itemize}
\end{enumerate}

\paragraph{Clasificación de Esquemas: Protocolos Sigma vs. Argumentos No Interactivos (NIZKAoK)}
Los sistemas de prueba de conocimiento cero se dividen estructuralmente según el modelo de interacción entre las partes:

\begin{itemize}
    \item \textbf{Protocolos Interactivos de Tres Movimientos ($\Sigma$-Protocolos):} Son la piedra angular de las construcciones algebraicas. Un protocolo $\Sigma$ opera en tres rondas consecutivas:
    \begin{enumerate}
        \item \textit{Compromiso (\textit{Commitment} $a$):} $\mathcal{P}$ genera un valor aleatorio interno $r$, calcula una promesa $a$ y la envía a $\mathcal{V}$.
        \item \textit{Desafío (\textit{Challenge} $e$):} $\mathcal{V}$ selecciona un desafío aleatorio $e \leftarrow \mathcal{E}$ de un espacio de desafíos y lo transmite a $\mathcal{P}$.
        \item \textit{Respuesta (\textit{Response} $z$):} $\mathcal{P}$ calcula la respuesta $z = f(x, w, r, e)$ y la envía a $\mathcal{V}$, quien evalúa una función de verificación pública $V(x, a, e, z) \in \{0, 1\}$ \cite{badiezadegan2026complexity}.
    \end{enumerate}
    Un protocolo $\Sigma$ satisface tradicionalmente la propiedad de \textit{Conocimiento Cero ante Verificador Honesto Especial} (\textit{Special HVZK}), lo que significa que el simulador puede generar transcripciones válidas $(a, e, z)$ eligiendo primero el desafío $e$ al azar y calculando $(a, z)$ de forma inversa.

    \item \textbf{Argumentos No Interactivos de Conocimiento (NIZKAoK) y Heurística de Fiat-Shamir:}
    En aplicaciones del mundo real, la interacción de múltiples rondas genera una latencia inaceptable. Para eliminar la interacción, se aplica la \textbf{Heurística de Fiat-Shamir} \cite{fiat1986how}. 
    
    En lugar de esperar el desafío aleatorio $e$ emitido por un verificador interactivo, el proponente computa el desafío internamente mediante el uso de una función hash criptográfica modelada como un \textbf{Oráculo Aleatorio} (\textit{Random Oracle Model} - ROM):
    \begin{equation}
    e = H(x \parallel a)
    \end{equation}
    La prueba no interactiva resultante es la tupla $\pi = (a, z)$. El verificador comprueba la validez calculando $e' = H(x \parallel a)$ y verificando que $V(x, a, e', z) = 1$. Esto transforma un protocolo $\Sigma$ interactivo en un Argumento No Interactivo de Conocimiento de Conocimiento Cero (\textit{Non-Interactive Zero-Knowledge Argument of Knowledge} - NIZKAoK) totalmente autónomo y verificable por terceros.
\end{itemize}

\paragraph{Ejemplos Didácticos e Intuitivos}
Para ilustrar intuitivamente los conceptos teóricos descritos, se presentan dos analogías clásicas y sus formulaciones matemáticas equivalentes:

\begin{itemize}
    \item \textbf{Ejemplo Interactivo: La Cueva de Alí Babá (Visualización de la Propiedad ZKP):}
    Imagine una cueva circular con una sola entrada/salida en la parte inferior y una puerta mágica en el extremo opuesto que requiere una palabra secreta (el testigo $w$) para abrirse. La puerta divide el camino en dos pasillos: la rama $A$ (izquierda) y la rama $B$ (derecha).
    \begin{itemize}
        \item \textit{Interacción:} El proponente $\mathcal{P}$ ingresa a la cueva y elige aleatoriamente caminar por el pasillo $A$ o $B$ sin que el verificador $\mathcal{V}$ lo vea. Luego, $\mathcal{V}$ se sitúa en la entrada y grita el desafío: ``¡Sal por el pasillo $A$!'' o ``¡Sal por el pasillo $B$!''.
        \item \textit{Completitud:} Si $\mathcal{P}$ conoce el secreto $w$, puede abrir la puerta mágica si es necesario y salir siempre por el pasillo solicitado por $\mathcal{V}$.
        \item \textit{Solidez:} Si un impostor $\mathcal{P}^*$ no conoce $w$, solo puede cumplir la orden si el pasillo pedido coincide con el pasillo en el que entró inicialmente (probabilidad $1/2$). Al repetir este proceso durante $k$ rondas independientes, la probabilidad de que el impostor engañe a $\mathcal{V}$ es de $(1/2)^k$. Para $k = 128$, la probabilidad es $2^{-128} \le \text{negl}(\lambda)$, garantizando solidez.
        \item \textit{Conocimiento Cero:} Un observador externo solo ve a $\mathcal{P}$ saliendo por el pasillo correcto. No aprende nada sobre la palabra secreta $w$. Además, un simulador podría fabricar una grabación de video idéntica grabando únicamente los intentos donde el pasillo solicitado coincida con la elección previa de $\mathcal{P}$, demostrando que la prueba no revela información sobre el secreto.
    \end{itemize}

    \item \textbf{Ejemplo No Interactivo: El Protocolo de Schnorr Convertido mediante Fiat-Shamir:}
    Considere el problema del logaritmo discreto sobre un grupo cíclico $\mathbb{Z}_q^*$ generado por $g$ de orden primo $q$ \cite{schnorr1991efficient, badiezadegan2026complexity}. La relación es $R_{DL} = \{ (y, x) \mid y = g^x \pmod q \}$.
    \begin{enumerate}
        \item \textit{Generación de la prueba no interactiva $\pi$:} $\mathcal{P}$ elige $r \leftarrow \mathbb{Z}_q$, calcula el compromiso $a = g^r \pmod q$, computa el desafío no interactivo $e = H(y \parallel a \parallel M)$ (donde $M$ puede ser un mensaje opcional) y calcula la respuesta $z = r + e \cdot x \pmod q$. La prueba generada es $\pi = (a, z)$.
        \item \textit{Verificación no interactiva:} Cualquier verificador recomputa el desafío $e = H(y \parallel a \parallel M)$ y comprueba si se satisface la ecuación:
        \begin{equation}
        g^z \stackrel{?}{=} a \cdot y^e \pmod q
        \end{equation}
        Dado que $g^z = g^{r + ex} = g^r \cdot (g^x)^e = a \cdot y^e \pmod q$, la verificación es matemáticamente exacta. El oráculo aleatorio $H$ actúa como una fuente impredecible de aleatoriedad, garantizando que el proponente no pueda alterar $a$ después de conocer $e$.
    \end{enumerate}
\end{itemize}

\subsubsection{Construcciones de ZKP sobre Retículos: De Stern a Lyubashevsky y el Framework LNP}

\paragraph{Evolución de los ZKP sobre Retículos: Paradigma Combinatorio vs. Fiat-Shamir con Abortos}
La transición de la criptografía basada en problemas algebraicos tradicionales (como la factorización de enteros o el logaritmo discreto) hacia la criptografía basada en retículos (\textit{lattice-based cryptography}) introdujo desafíos estructurales significativos en el diseño de pruebas de conocimiento cero.

Históricamente, el primer protocolo fundamental de ZKP sobre retículos fue el esquema de Stern (1993), concebido originalmente para el problema de decodificación por síndrome (\textit{Syndrome Decoding}) y posteriormente adaptado al problema \textit{Short Integer Solution} (SIS). El protocolo de Stern se basa en un enfoque combinatorio de permutaciones aleatorias y compromisos ocultos. Sin embargo, su limitación principal reside en que posee un alto error de solidez por ronda de $2/3$. Para alcanzar un nivel de seguridad adecuado de $128$ bits, el protocolo requiere repetirse más de 219 veces ($\approx (2/3)^{219} \le 2^{-128}$), lo que genera pruebas extremadamente pesadas con un consumo de ancho de banda de varios cientos de kilobytes \cite{badiezadegan2026complexity, cayrel2012improved}.

Para superar estas ineficiencias combinatorias, Cayrel et al. \cite{cayrel2012improved} desarrollaron una construcción mejorada basada en el problema SIS sobre retículos modulares. Mediante el uso de funciones de compromiso ($COM$) estadísticamente ocultas y computacionalmente vinculantes, combinadas con vectores de máscara aleatorios $u \in \mathbb{Z}_q^m$ y matrices de permutación $P_\sigma \in S_m$, Cayrel et al. lograron reducir el error de solidez por ronda a:
\begin{equation}
\epsilon = \frac{q+1}{2q} \approx \frac{1}{2}
\end{equation}
Esta reducción de $2/3$ a $1/2$ disminuyó drásticamente el número de repeticiones requeridas para alcanzar la solidez post-cuántica deseada, reduciendo la sobrecarga de almacenamiento y comunicación en retículos ideales (\textit{Ideal-SIS}) \cite{cayrel2012improved}.

No obstante, el cambio de paradigma decisivo ocurrió con la introducción de la técnica de \textbf{Fiat-Shamir con Abortos} (\textit{Fiat-Shamir with aborts}) ideada por Lyubashevsky \cite{lyubashevsky2009fiat, lyubashevsky2012lattice}. En los protocolos $\Sigma$ tradicionales sobre grupos (como Schnorr), la respuesta del proponente adopta la forma $z = r + e \cdot x$. Si se intentara aplicar esta fórmula directamente sobre un secreto vectorial corto $s$ en un retículo, el valor de $z = y + c \cdot s$ revelaría información parcial sobre la clave secreta $s$ a través de la norma y distribución de $z$. 

Para evitar la fuga de información sin necesidad de ocultar la estructura mediante permutaciones pesadas, Lyubashevsky introdujo la técnica de \textbf{muestreo por rechazo} (\textit{rejection sampling}): el proponente calcula $z = y + c \cdot s$ y comprueba si la distribución de $z$ depende del secreto $s$. Si $z$ excede un determinado umbral de norma o se desvía de una distribución uniforme/gaussiana independiente de $s$, el proponente \textbf{aborta} la ejecución de la ronda y recomienza el proceso con una nueva aleatoriedad $y$. Esto garantiza rigurosamente la propiedad de cero conocimiento mientras genera respuestas algebraicas directas y compactas \cite{lyubashevsky2012lattice}.

\paragraph{Formulación de Relaciones Lineales y Restricciones Cuadráticas sobre Retículos}
Los sistemas modernos de ZKP sobre retículos estructuran las afirmaciones criptográficas mediante relaciones algebraicas modulares definidas sobre un anillo o cuerpo finito $\mathbb{Z}_q$.

\begin{enumerate}
    \item \textbf{Relación Inhomogeneous Small Integer Solution (ISIS):}
    Dado un módulo $q \in \mathbb{N}$, una matriz pública $A \in \mathbb{Z}_q^{n \times m}$, un vector de síndrome público $u \in \mathbb{Z}_q^n$ y una cota de norma $\mathcal{B}_s > 0$, la relación de conocimiento para el problema ISIS se define formalmente como:
    \begin{equation}
    R_{ISIS} = \left\{ \left( (A, u, \mathcal{B}_s), s \right) \;\middle|\; A s \equiv u \pmod q \;\;\land\;\; \|s\| \le \mathcal{B}_s \right\}
    \end{equation}
    Akleylek y Soysaldı \cite{akleylek2021post} formalizan cómo esta relación permite construir esquemas de autenticación y ZKP post-cuánticos eficientes para entornos restringidos como IoT. La seguridad del protocolo descansa en la dureza de encontrar un vector no nulo corto $s$ que pertenezca al conjunto de soluciones del sistema lineal en el peor de los casos sobre retículos modulares $L_q^\perp(A) = \{x \in \mathbb{Z}^m \mid Ax \equiv 0 \pmod q\}$ \cite{akleylek2021post}.

    \item \textbf{Relaciones para Firmas e Identidades ($R_{Sign}$) y Restricciones Cuadráticas:}
    En protocolos de acreditación de identidad, firmas digitales y pruebas de privacidad avanzadas, no basta con probar relaciones puramente lineales $As = u \pmod q$. Se requiere demostrar que el testigo corto $s$ satisface simultáneamente un conjunto de restricciones cuadráticas modulares y ecuaciones bilineales:
    \begin{equation}
    R_{Quad} = \left\{ \left( (A, \{B_i\}, c), s \right) \;\middle|\; A s + \sum_{i} s_i B_i s \equiv c \pmod q \;\;\land\;\; \|s\|_\infty \le \mathcal{B}_s \right\}
    \end{equation}
    Estas relaciones cuadráticas permiten verificar en conocimiento cero que una clave secreta está correctamente firmada por una autoridad emisora sin revelar el contenido de la credencial ni la firma misma.
\end{enumerate}

\paragraph{El Framework LNP (Lyubashevsky-Nguyen-Plançon) y Construcciones de Pruebas Cortas}
Para unificar y optimizar la demostración de relaciones lineales y cuadráticas sobre retículos, Lyubashevsky, Nguyen y Plançon desarrollaron el \textbf{Framework LNP} \cite{lyubashevsky2019framework}, extendido posteriormente por Bootle et al. \cite{baum2020concretely}.

El Framework LNP opera sobre retículos en módulos definidos en el anillo de polinomios $R_q = \mathbb{Z}_q[X]/(X^n + 1)$ y se fundamenta en los siguientes componentes algebraicos:
\begin{itemize}
    \item \textbf{Esquema de Compromiso BDLOP:} Un esquema de compromiso de tipo Pedersen adaptado a retículos que es aditivamente homomórfico y computacionalmente vinculante bajo la dureza del problema Module-SIS (M-SIS) y estadísticamente oculto bajo Module-LWE (M-LWE) \cite{lyubashevsky2019framework}.
    \item \textbf{Pruebas de Producto Escalar en Retículos (\textit{Lattice Inner-Product Proofs}):} Permite al proponente convencer al verificador de que un conjunto de vectores comprometidos satisface ecuaciones cuadráticas de la forma $\langle s_1, s_2 \rangle = c$, reduciendo el tamaño de la prueba de manera equivalente a las técnicas tipo \textit{Bulletproofs}, pero basadas exclusivamente en la dureza de retículos.
    \item \textbf{Integración con Abortos:} Combina el protocolo $\Sigma$ no interactivo (vía Fiat-Shamir) con parámetros de rechazo optimizados, garantizando que el tamaño total de la prueba se mantenga en el orden de pocos kilobytes (típicamente entre $8 \text{ KB}$ y $15 \text{ KB}$), en contraste con los megabytes requeridos por los esquemas combinatorios tipo Stern \cite{lyubashevsky2019framework, ma2026laasso}.
\end{itemize}

\subsubsection{El Obstáculo del Rebobinado Cuántico (Quantum Rewinding Barrier)}

\paragraph{Mecanismo Clásico de Rebobinado y el Lemma de Bifurcación (\textit{Forking Lemma})}
En la criptografía clásica, las demostraciones formales de solidez (\textit{soundness}) y la propiedad de argumento de conocimiento para protocolos $\Sigma$ e identificaciones no interactivas se basan fundamentalmente en la técnica del \textbf{rebobinado} (\textit{rewinding}) y en el \textit{Forking Lemma} (introducido por Pointcheval y Stern).

Bajo este paradigma, la prueba de seguridad se construye mediante un \textbf{Extractor de Conocimiento} ($\mathcal{E}$) determinista o probabilístico en tiempo polinomial que interactúa con un proponente deshonesto $\mathcal{P}^*$ como una ``caja negra'':
\begin{enumerate}
    \item $\mathcal{E}$ ejecuta a $\mathcal{P}^*$ hasta que este emite un compromiso inicial $a$.
    \item $\mathcal{E}$ envía un primer desafío aleatorio $e_1$ y registra la respuesta válida $z_1$, obteniendo la transcripción $(a, e_1, z_1)$.
    \item A continuación, $\mathcal{E}$ \textbf{rebobina} el estado de ejecución interno de $\mathcal{P}^*$ al punto exacto inmediatamente posterior al envío del compromiso $a$.
    \item En este estado restaurado, $\mathcal{E}$ emite un segundo desafío diferente $e_2 \neq e_1$ y obtiene una segunda respuesta válida $z_2$, generando la transcripción $(a, e_2, z_2)$.
\end{enumerate}

Al disponer de dos transcripciones válidas $((a, e_1, z_1), (a, e_2, z_2))$ que comparten el mismo compromiso $a$, el algoritmo extractor $\mathcal{E}$ aplica relaciones algebraicas directas para despejar y extraer el testigo secreto $w$ en tiempo polinomial, reduciendo la solidez de la prueba a la dureza del problema matemático subyacente.

\paragraph{Falla Fundamental en el Modelo Cuántico: Colapso de la Función de Onda y No-Clonación}
Cuando el adversario es una entidad cuántica, el mecanismo clásico de rebobinado colapsa categóricamente. Lombardi et al. \cite{barak2003lower, liu2019revisiting} formalizaron las restricciones matemáticas que impiden adaptar directamente el \textit{rewinding} al entorno post-cuántico.

Un proponente deshonesto cuántico $\mathcal{P}^*$ almacena su estado interno como un estado cuántico complejo en superposición, representado por un vector de estado $|\psi\rangle$ en un espacio de Hilbert:
\begin{equation}
|\psi\rangle = \sum_{i} \alpha_i |i\rangle, \quad \text{con } \sum_{i} |\alpha_i|^2 = 1
\end{equation}
Cuando el extractor cuántico $\mathcal{E}$ interactúa con $\mathcal{P}^*$ para obtener la respuesta $z_1$ al primer desafío $e_1$, efectúa inevitables mediciones cuánticas sobre los registros de $\mathcal{P}^*$. Por el principio de medición von Neumann:
\begin{itemize}
    \item \textbf{Colapso del Estado:} La medición provoca el colapso irreversible de la función de onda de $|\psi\rangle$ a un autostado proyectado $|\psi_{e_1, z_1}\rangle$. Trazas o fases de la superposición previa son destruidas por decoherencia y medición.
    \item \textbf{Imposibilidad de Clonación:} Por el \textbf{Teorema de No-Clonación} (\textit{No-Cloning Theorem}), el extractor $\mathcal{E}$ no puede realizar una copia de respaldo del estado cuántico desconocido $|\psi\rangle$ antes de enviar el primer desafío $e_1$.
\end{itemize}

En consecuencia, el extractor es incapaz de ``rebobinar'' o restaurar a $\mathcal{P}^*$ a su estado cuántico exacto $\vert{}\psi\rangle$ previo a la primera medición. Si $\mathcal{E}$ intenta enviar un segundo desafío $e_2$ sobre el estado alterado $\vert{}\psi_{e_1, z_1}\rangle$, la respuesta obtenida no guardará la coherencia necesaria con el compromiso $a$, invalidando el proceso de extracción del testigo \cite{barak2003lower}.

\paragraph{Insuficiencia de las Reducciones Clásicas para Seguridad Poscuántica}
La existencia del obstáculo del rebobinado cuántico conduce a un resultado crítico en la criptografía post-cuántica: \textbf{la dureza post-cuántica del problema matemático subyacente (como LWE o SIS) no garantiza automáticamente la seguridad post-cuántica del protocolo ZKP si su demostración de solidez depende de reducciones por rebobinado} \cite{liu2019revisiting}.

Incluso si un problema como \textit{Learning With Errors} (LWE) es resistente a algoritmos cuánticos de tiempo polinomial (como el algoritmo de Shor), un atacante cuántico podría violar la propiedad de solidez o de cero conocimiento explotando la incapacidad del extractor para ejecutar el \textit{rewinding} sin destruir el estado de superposición. Por este motivo, el diseño de argumentos no interactivos de conocimiento poscuánticos (NIZKAoK) exige modelos analíticos más rigurosos, tales como:
\begin{enumerate}
    \item El \textbf{Modelo de Oráculo Aleatorio Cuántico} (\textit{Quantum Random Oracle Model} - QROM), donde el adversario cuántico puede realizar consultas en superposición a las funciones hash.
    \item La propiedad de \textbf{Extractabilidad en Línea Recta} (\textit{Straight-Line Extractability} - SLE), la cual permite la extracción del testigo en una única ejecución sin necesidad de rebobinado.
\end{enumerate}

\subsubsection{Pruebas NIZKAoK con Extractabilidad en Línea Recta (Straight-Line Extractability - SLE)}

\paragraph{Definición Formal de NIZKAoK con Extractabilidad en Línea Recta y Multi-Prueba}
Para solventar de manera definitiva la incompatibilidad matemática entre el rebobinado clásico y la mecánica cuántica, la criptografía poscuántica adopta el paradigma de \textbf{Extractabilidad en Línea Recta} (\textit{Straight-Line Extractability} - SLE). En un sistema de Argumento No Interactivo de Conocimiento de Conocimiento Cero (NIZKAoK) con propiedad SLE, la extracción del testigo secreto $w$ se realiza en una sola pasada de ejecución directa (\textit{single forward pass}), anulando la necesidad de restaurar estados o repetir desafíos \cite{barak2003lower, ma2026laasso}.

Un esquema NIZKAoK con extractabilidad en línea recta sobre una relación NP $R = \{(x, w)\}$ consta de un triplete de algoritmos probabilísticos en tiempo polinomial $\Pi = (\text{Setup}, \mathcal{P}, \mathcal{V})$ junto con un par Extractor $\mathcal{E} = (\mathcal{E}_1, \mathcal{E}_2)$:
\begin{enumerate}
    \item \textbf{Generación de Parámetros con Trampilla:} $\mathcal{E}_1(1^\lambda) \to (\sigma, \tau)$. El algoritmo extractor genera una Cadena de Referencia Común (CRS) $\sigma$ computacionalmente indistinguible de una CRS generada honestamente por $\text{Setup}(1^\lambda)$, junto con una información secreta de trampilla (\textit{trapdoor}) $\tau$.
    \item \textbf{Extracción Directa:} $\mathcal{E}_2(\sigma, \tau, x, \pi) \to w'$. Dado el par $(\sigma, \tau)$, una instancia $x$ y una prueba no interactiva válida $\pi$, $\mathcal{E}_2$ extrae directamente un testigo $w'$ tal que $(x, w') \in R$.
\end{enumerate}

La propiedad de \textbf{Extractabilidad Multi-Prueba} (\textit{Multi-Proof Extractability}) extiende formalmente esta noción al entorno donde un adversario deshonesto emite una secuencia adaptativa de múltiples demostraciones $\{\pi_1, \pi_2, \dots, \pi_k\}$ sobre diferentes instancias $\{x_1, x_2, \dots, x_k\}$. La trampilla estática $\tau$ incrustada en la CRS permite al extractor $\mathcal{E}_2$ derivar individualmente cada testigo $w_i$ en tiempo polinomial sin alterar en ningún momento el canal de comunicación ni interferir con el estado cuántico del proponente \cite{ma2026laasso}.

\paragraph{Mecanismo de Operación del Extractor: Trampas en la CRS y el Modelo QROM}
La arquitectura de un extractor en línea recta opera abstrayendo la verificación del canal interactivo mediante el uso de esquemas de compromiso en modo dual (\textit{dual-mode commitments}) o funciones de cifrado con trampilla de pérdida (\textit{lossy encryption}) \cite{barak2003lower, ma2026laasso}:

\begin{itemize}
    \item \textbf{Incrustación de Trampillas en la CRS:} En el modo de extracción, la CRS publica una clave de compromiso o clave pública de reticulado que contiene una trampilla algebraica $\tau$ (por ejemplo, una base corta de un retículo $L_q^\perp(\mathbf{A})$ bajo el problema LWE/SIS). Durante la construcción de la prueba no interactiva $\pi = (\mathbf{a}, \mathbf{z})$, el proponente emite un compromiso o cifrado del testigo $\mathbf{w}$. Mientras que para un verificador ordinario la CRS garantiza el ocultamiento perfecto de $\mathbf{w}$, el extractor $\mathcal{E}_2$, provisto de $\tau$, utiliza la trampilla para resolver el problema algebraicoy descifrar/descomprometer el testigo de manera determinista y directa.

    \item \textbf{Extracción en el Modelo de Oráculo Aleatorio Cuántico (QROM):} Cuando la prueba no interactiva se construye en el QROM, los atacantes cuánticos pueden realizar consultas a la función hash $H$ en superposición de estados $\sum_{x, y} \alpha_{x,y} \vert{}x, y\rangle$. Para lograr la extractabilidad en este modelo sin rebobinado, esquemas como la transformación de Unruh o las variantes optimizadas del Framework LNP incrustan la respuesta en la propia estructura del oráculo reprogramable de forma determinista, permitiendo que la lectura del registro por parte de $\mathcal{E}_2$ recupere $w$ sin colapsar el estado del adversario \cite{liu2019revisiting, ma2026laasso}.
\end{itemize}

\paragraph{Demostración de la Solidez Poscuántica Estricta}
La solidez poscuántica estricta de un sistema NIZKAoK basado en SLE se formaliza mediante la siguiente reducción:

Sea $\mathcal{P}^*$ un proponente cuántico en tiempo polinomial (QPT) capaz de generar una prueba no interactiva válida $\pi$ para una instancia $x \notin L_R$ con probabilidad no despreciable $\epsilon(\lambda)$. Se define el experimento de extracción sobre el par de algoritmos $(\mathcal{E}_1, \mathcal{E}_2)$:
\begin{equation}
\Pr \left[ 
\begin{array}{c} 
(\sigma, \tau) \leftarrow \mathcal{E}_1(1^\lambda), \\ 
(x, \pi) \leftarrow \mathcal{P}^*(\sigma), \\ 
w' \leftarrow \mathcal{E}_2(\sigma, \tau, x, \pi) 
\end{array} 
\;\middle|\; 
\mathcal{V}(\sigma, x, \pi) = 1 \;\land\; (x, w') \notin R 
\right] \le \text{negl}(\lambda)
\end{equation}

Puesto que el extractor $\mathcal{E}_2$ opera en una sola línea recta (\textit{straight-line}) sin ejecutar interrupciones, mediciones sobre registros de superposición interna, ni peticiones de reinicio al sistema, no existe perturbación sobre la función de onda del adversario cuántico. En consecuencia, la solidez del protocolo es \textbf{absoluta y strictly poscuántica}, garantizando que ningún algoritmo cuántico pueda falsificar una prueba sin que el testigo equivalente sea recuperable instantáneamente mediante la trampilla de la CRS \cite{barak2003lower, ma2026laasso}.

\subsubsection{Privacidad del Canal y Verificadores Designados: El Mecanismo de Token Hiding}

\paragraph{El Problema del Verificador No Designado y la Intercepción por Terminales Comprometidos}
En los protocolos de autenticación distribuidos y sistemas de inicio de sesión único (\textit{Single Sign-On} - SSO), las credenciales y pruebas criptográficas generadas por el usuario deben transmitirse frecuentemente a través de intermediarios no confiables o redes públicas antes de alcanzar la entidad verificadora final \cite{ma2026laasso}.

Si las pruebas de conocimiento cero (ZKP) o los tokens de autenticación se transmitieran en texto claro a través del canal de comunicación, surgiría el problema del \textit{verificador no designado} \cite{ma2026laasso}. Un cajero automático comprometido mediante soporte malicioso o dispositivos de captura ilegítima (\textit{skimmers}) podría inspeccionar el contenido del token, extraer metadatos de la transacción o rastrear las pruebas enviadas en sesiones sucesivas \cite{ma2026laasso, liu2019revisiting}. 

Para prevenir estos ataques de rastreabilidad y proteger la identidad del usuario frente a intermediarios deshonestos, los sistemas modernos de NIZKAoK integran el mecanismo de \textbf{ocultamiento de tokens} (\textit{token hiding}), garantizando que la prueba generada solo sea descifrable y verificable por el servidor bancario legítimo (verificador designado).

\subsubsection{Privacidad del Canal y Verificadores Designados: El Mecanismo de Token Hiding}

\paragraph{El Problema del Verificador No Designado y la Intercepción por Terminales Comprometidos}
En los protocolos de autenticación distribuidos y sistemas de inicio de sesión único (\textit{Single Sign-On} - SSO), las credenciales y pruebas criptográficas generadas por el usuario deben transmitirse frecuentemente a través de intermediarios no confiables o redes públicas antes de alcanzar la entidad verificadora final \cite{ma2026laasso}.

Si las pruebas de conocimiento cero (ZKP) o los tokens de autenticación se transmitieran en texto claro a través del canal de comunicación, surgiría el problema del \textit{verificador no designado} \cite{ma2026laasso}. Un cajero automático comprometido mediante soporte malicioso o dispositivos de captura ilegítima (\textit{skimmers}) podría inspeccionar el contenido del token, extraer metadatos de la transacción o rastrear las pruebas enviadas en sesiones sucesivas \cite{ma2026laasso, liu2019revisiting}. 

Para prevenir estos ataques de rastreabilidad y proteger la identidad del usuario frente a intermediarios deshonestos, los sistemas modernos de NIZKAoK integran el mecanismo de \textbf{ocultamiento de tokens} (\textit{token hiding}), garantizando que la prueba generada solo sea descifrable y verificable por el servidor bancario legítimo (verificador designado).

\subsubsection{Propiedades Avanzadas de Privacidad: No Vinculabilidad y Denegabilidad Total}

\paragraph{No Vinculabilidad de Credenciales y Tokens (\textit{Unlinkability})}
En protocolos de autenticación anonimizados y sistemas transaccionales poscuánticos, la \textbf{no vinculabilidad} (\textit{unlinkability}) garantiza que un adversario o servidor intermedio no pueda determinar si dos o más ejecuciones del protocolo provienen del mismo usuario o dispositivo \cite{ma2026laasso}.

Formalmente, consideremos dos sesiones de autenticación $i$ y $j$ realizadas por el mismo cliente utilizando una clave secreta $s$, las cuales producen los tokens de autenticación $\text{tok}^{(i)}$ y $\text{tok}^{(j)}$. La propiedad de no vinculabilidad establece que para cualquier adversario en tiempo polinomial probabilístico (PPT) $\mathcal{A}$, la ventaja de distinguir entre dos tokens generados por el mismo usuario y dos tokens generados por usuarios totalmente distintos es infinitesimal \cite{ma2026laasso}:
\begin{equation}
\left| \Pr\left[ \mathcal{A}\left(\text{tok}^{(i)}(s), \text{tok}^{(j)}(s)\right) = 1 \right] - \Pr\left[ \mathcal{A}\left(\text{tok}^{(i)}(s_1), \text{tok}^{(j)}(s_2)\right) = 1 \right] \right| \le \text{negl}(\lambda)
\end{equation}

En la arquitectura propuesta, esta garantía se logra mediante dos capas de aleatorización:
\begin{enumerate}
    \item \textbf{Aleatorización de Máscaras sobre el Retículo:} En cada iteración del argumento NIZKAoK, la respuesta $z = y + c \cdot s$ y el compromiso $a$ emplean vectores de máscara efímeros $y \stackrel{\$}{\leftarrow} D_{\sigma}^m$ totalmente independientes tomados de una distribución gaussiana discreta (o uniforme), lo que invalida cualquier correlación algebraica estática entre transacciones \cite{lyubashevsky2012lattice, ma2026laasso}.
    \item \textbf{Cifrado Híbrido con Claves Efímeras:} Tal como se demostró en el mecanismo de \textit{token hiding}, cada token $\text{tok} = (\text{tok}_1, \text{tok}_2)$ se encapsula mediante una clave simétrica efímera $K$ de un solo uso bajo Kyber (ML-KEM) y AES-256 \cite{ma2026laasso}. Puesto que ambos componentes resultan indistinguibles de cadenas aleatorias uniformes sobre $\{0,1\}^*$, dos transacciones del mismo usuario exhiben distribuciones numéricas estadísticamente independientes y no vinculables.
\end{enumerate}

\paragraph{Denegabilidad Total (\textit{Full Deniability}) y Simulación de Sesión}
La propiedad de \textbf{denegabilidad total} (\textit{full deniability} / \textit{plausible deniability}), analizada formalmente por An et al. \cite{an2023secrethandshakes}, resulta indispensable para proteger al usuario contra coerciones o auditorías maliciosas posteriores a la transacción (\textit{post-hoc coercion}). Si un atacante o una entidad reguladora intercepta la transcripción de comunicación completa de una sesión pasada $\mathcal{T}$, la denegabilidad total asegura que el usuario puede negar legítimamente haber participado en dicha interacción, demostrando que cualquier tercero o el propio verificador podría haber fabricado la transcripción por sí mismo \cite{an2023secrethandshakes}.

Matemáticamente, la denegabilidad total se fundamenta en la existencia de un \textbf{Simulador de Transcripción} $\mathcal{S}_{den}$ que puede ejecutarse utilizando únicamente la clave pública del verificador o parámetros públicos de la CRS, sin requerir la clave privada del usuario $s$ ni la cooperación activa de este \cite{an2023secrethandshakes}:
\begin{equation}
\mathcal{T}_{real} = \langle \mathcal{P}(s), \mathcal{V}(vpk) \rangle \;\approx_c\; \mathcal{S}_{den}(vpk) = \mathcal{T}_{sim}
\end{equation}

En un protocolo con verificador designado y denegabilidad total (como los esquemas de \textit{secret handshakes} de An et al. \cite{an2023secrethandshakes}), la clave de verificación o la trampa del verificador le confiere a este la capacidad técnica de simular transcripciones indistinguibles de las reales. De este modo, si la transcripción $\mathcal{T}$ se expone a un tribunal u observador externo, carece de valor probatorio legal o criptográfico: el verificador no puede probar a un tercero que la prueba fue generada por el usuario y no sintetizada por el propio verificador \cite{an2023secrethandshakes, ma2026laasso}.

\subsubsection{Estado del Arte Comparativo de Frameworks ZKP Poscuánticos}

\paragraph{Análisis Comparativo de Paradigmas Poscuánticos y Resistencia a Rebobinado Cuántico}
La evolución de las pruebas de conocimiento cero en la era poscuántica está definida por la capacidad de los protocolos para mantener la solidez frente a adversarios con capacidades de computación cuántica. Mientras que los esquemas pioneros dependían exclusivamente de análisis de seguridad clásicos basados en el \textit{Forking Lemma}, las construcciones contemporáneas han adoptado modelos adaptados al Modelo de Oráculo Aleatorio Cuántico (QROM) y técnicas de Extractabilidad en Línea Recta (\textit{Straight-Line Extractability} - SLE) \cite{barak2003lower, liu2019revisiting}.

\begin{table}[h!]
\centering
\small
\caption{Comparativa Teórica de Paradigmas ZKP frente a Atacantes Cuánticos y Extracción}
\label{tab:zkp_quantum_comparison}
\begin{tabular}{|p{3.2cm}|p{2.8cm}|p{3.2cm}|p{2.2cm}|p{3.2cm}|}
\hline
\textbf{Paradigma / Esquema} & \textbf{Problema Subyacente} & \textbf{Mecanismo de Extracción} & \textbf{Modelo de Seguridad} & \textbf{Garantía Poscuántica (QROM)} \\ \hline
Stern (1993) / Cayrel et al. (2012) \cite{cayrel2012improved} & SIS / Decodificación por Síndrome & Rebobinado Clásico (\textit{Forking}) & Clásico (ROM) & \textbf{Vulnerable} al colapso de la función de onda \cite{barak2003lower}. \\ \hline
Fiat-Shamir con Abortos (Lyubashevsky) \cite{lyubashevsky2012lattice} & SIS / LWE sobre Retículos & Rebobinado Clásico / Muestreo por Rechazo & QROM & \textbf{Condicional} (Requiere parámetros de aborto expandidos) \cite{liu2019revisiting}. \\ \hline
Framework LNP / Bootle et al. \cite{lyubashevsky2019framework, baum2020concretely} & M-SIS / M-LWE (Anillos $R_q$) & Rebobinado Adaptado / BDLOP & QROM & \textbf{Robusto} mediante reducciones strictly sobre retículos \cite{lyubashevsky2019framework}. \\ \hline
\textbf{Arquitectura LAASSO (SLE)} \cite{ma2026laasso} & Module-LWE / Module-SIS & \textbf{Extractabilidad en Línea Recta (SLE)} & QROM + Trampilla CRS & \textbf{Estrictamente Poscuántica} (Inmune a fallos de rebobinado) \cite{barak2003lower, ma2026laasso}. \\ \hline
\end{tabular}
\end{table}

\paragraph{Evaluación Cuantitativa de Métricas de Rendimiento}
La viabilidad operativa de un sistema ZKP poscuántico en entornos reales (tales como dispositivos de borde, tarjetas inteligentes o transacciones financieras distribuidas) se mide a través del tamaño de la prueba generada, la latencia computacional de generación/verificación y la complejidad de las operaciones en el anillo $\mathbb{Z}_q$ o $R_q = \mathbb{Z}_q[X]/(X^n + 1)$ \cite{badiezadegan2026complexity, lyubashevsky2019framework, ma2026laasso}.

Las construcciones combinatorias tipo Stern sufren de una sobrecarga de tamaño inaceptable para canales de bajo ancho de banda, requiriendo cientos de kilobytes debido a las múltiples repeticiones necesarias para mitigar el error de solidez de $2/3$ \cite{badiezadegan2026complexity, cayrel2012improved}. En contraste, los marcos basados en polinomios modulares e integraciones de producto escalar (como LNP y LAASSO) reducen el tamaño de la prueba a un rango de $4 \text{ a } 15 \text{ KB}$ mediante multiplicaciones polinómicas eficientes ejecutadas con la Transformada de Número Entero (\textit{Number Theoretic Transform} - NTT) \cite{lyubashevsky2019framework, ma2026laasso}.

\begin{table}[h!]
\centering
\small
\caption{Matriz Técnico-Cuantitativa de Frameworks ZKP Poscuánticos}
\label{tab:zkp_metrics_matrix}
\begin{tabular}{|p{3.4cm}|p{2.3cm}|p{2.2cm}|p{3.5cm}|p{3.2cm}|}
\hline
\textbf{Framework / Protocolo} & \textbf{Tamaño de Prueba (KB)} & \textbf{Tiempo Gen. ($\text{ms}$)} & \textbf{Operaciones Dominantes ($\mathbb{Z}_q / R_q$)} & \textbf{Relación Soportada} \\ \hline
Pruebas Combinatorias (Stern / Cayrel) \cite{cayrel2012improved} & $100 - 300 \text{ KB}$ & $150 - 500 \text{ ms}$ & Permutaciones $P_\sigma$ y funciones hash & Relaciones SIS binarias/trinarias. \\ \hline
Circuitos Aritméticos Generales (Baum et al.) \cite{baum2020concretely} & $20 - 50 \text{ KB}$ & $50 - 150 \text{ ms}$ & Productos internos modulares sobre $R_q$ & Circuitos aritméticos arbitrarios. \\ \hline
Framework LNP (Lyubashevsky et al.) \cite{lyubashevsky2019framework} & $8 - 15 \text{ KB}$ & $10 - 30 \text{ ms}$ & Multiplicación Matriz-Vector NTT en $R_q$ & Restricciones lineales y cuadráticas. \\ \hline
\textbf{Arquitectura LAASSO (SLE + Token Hiding)} \cite{ma2026laasso} & \textbf{$4 - 9 \text{ KB}$} & \textbf{$2 - 10 \text{ ms}$} & Multiplicación NTT + Encapsulamiento Kyber/AES & $R_{ISIS}$, $R_{Sign}$ y ocultación de tokens. \\ \hline
\end{tabular}
\end{table}

\paragraph{Síntesis Teórica y Selección de la Arquitectura}
La integración de las propiedades analizadas determina que la arquitectura compuesta por \textbf{NIZKAoK con Extractabilidad en Línea Recta (SLE)} y \textbf{Ocultación de Tokens (\textit{Token Hiding}) basada en Kyber/AES} \cite{ma2026laasso} representa la solución óptima para el prototipo del marco teórico. Esta combinación garantiza:
\begin{enumerate}
    \item \textbf{Inmunidad Cuántica Absoluta:} Al prescindir del rebobinado, elimina la vulnerabilidad derivada del colapso del estado cuántico demostrada por Lombardi et al. \cite{barak2003lower}.
    \item \textbf{Eficiencia Extrema:} Ofrece el menor tamaño de prueba ($4 - 9 \text{ KB}$) y latencia ($2 - 10 \text{ ms}$), permitiendo su ejecución directa en chips de tarjetas o dispositivos IoT \cite{akleylek2021post, ma2026laasso}.
    \item \textbf{Privacidad del Canal:} Transforma el canal de transmisión en un mediador ciego mediante el enmascaramiento con Kyber (ML-KEM), previniendo ataques de seguimiento, \textit{skimming} o perfilamiento de usuarios en nodos intermedios \cite{ma2026laasso}.
\end{enumerate}

\subsection{Learning With Errors (LWE)}

\subsubsection{Introducción a la Criptografía Basada en Retículos}

\paragraph{Definición Matemática de Retículo y Geometría Discreta}

En el ámbito del álgebra lineal y la geometría discreta, un retículo (o \textit{lattice}) $\Lambda$ en el espacio euclídeo $m$-dimensional $\mathbb{R}^m$ se define formalmente como un subgrupo discreto del espacio vectorial real generado por todas las combinaciones lineales enteras de $k$ vectores linealmente independientes $\mathbf{b}_1, \mathbf{b}_2, \dots, \mathbf{b}_k \in \mathbb{R}^m$ (con $m \ge k$) \cite{huang2018lattice, ma2026laasso}:

\begin{equation}
\Lambda = \mathcal{L}(\mathbf{B}) = \left\{ \sum_{i=1}^{k} c_i \mathbf{b}_i \;\middle|\; c_i \in \mathbb{Z}, \, \mathbf{b}_1, \dots, \mathbf{b}_k \in \mathbb{R}^m \right\}
\end{equation}

donde $\mathbf{B} = [\mathbf{b}_1 | \mathbf{b}_2 | \dots | \mathbf{b}_k] \in \mathbb{R}^{m \times k}$ representa la matriz de base del retículo, mientras que los enteros $k$ y $m$ denotan el rango y la dimensión del retículo, respectivamente \cite{huang2018lattice}. Cuando el rango iguala a la dimensión ($k = m$), se afirma que el retículo es de rango completo (\textit{full-rank}) \cite{ma2026laasso}. Geométricamente, un retículo constituye una estructura periódica discreta e infinita sobre el espacio vectorial \cite{huang2018lattice}.

En el diseño de protocolos criptográficos poscuánticos, adquieren especial relevancia los retículos modulares discretos definidos sobre el anillo de enteros $\mathbb{Z}_q = \mathbb{Z}/q\mathbb{Z}$ (donde $q$ es un módulo primo o entero), conocidos como retículos $q$-arios (\textit{$q$-ary lattices}) \cite{huang2018lattice, an2023secrethandshakes}. Para una matriz modular $\mathbf{A} \in \mathbb{Z}_q^{n \times m}$, se definen los siguientes conjuntos duales fundamentales:

\begin{align}
\Lambda_q^\perp(\mathbf{A}) &= \left\{ \mathbf{y} \in \mathbb{Z}^m \;\middle|\; \mathbf{A}\mathbf{y} \equiv \mathbf{0} \pmod q \right\} \label{eq:q_ary_perp} \\
\Lambda_q(\mathbf{A}) &= \left\{ \mathbf{y} \in \mathbb{Z}^m \;\middle|\; \mathbf{y} \equiv \mathbf{A}^T\mathbf{s} \pmod q, \, \text{para algún } \mathbf{s} \in \mathbb{Z}_q^n \right\}
\end{align}

así como el retículo desplazado $\Lambda_u(\mathbf{A}) = \{ \mathbf{x} \in \mathbb{Z}^m \mid \mathbf{A}\mathbf{x} \equiv \mathbf{u} \pmod q \}$ para un vector $u \in \mathbb{Z}_q^n$ \cite{an2023secrethandshakes}. Asimismo, sobre anillos ciclotómicos $\mathcal{R}_q = \mathbb{Z}_q[X]/\langle X^d + 1 \rangle$, la estructura puede extenderse a retículos ideales o de tipo NTRU $\Lambda_{h,q} = \{ (u,v) \in \mathcal{R}_q \times \mathcal{R}_q \mid u + vh \equiv 0 \pmod q \}$ \cite{ma2026laasso}.

\paragraph{Problemas Duros del Peor Caso en Retículos}

La solidez de la criptografía basada en retículos reposa sobre la intratabilidad computacional de diversos problemas geométricos multidimensionales \cite{huang2018lattice, rewal2023quantumsafe, ma2026laasso, an2023secrethandshakes}:

\begin{itemize}
    \item \textbf{Problema del Vector más Corto (Shortest Vector Problem - SVP):} Dada una base arbitraria $\mathbf{B}$ de un retículo $L$, el objetivo consiste en hallar el vector no nulo de menor norma euclídea posible en $L$ \cite{huang2018lattice}.
    
    \item \textbf{Problema del Vector más Cercano (Closest Vector Problem - CVP):} Dada una base $\mathbf{B}$ de un retículo $L$ y un vector objetivo $\mathbf{t} \notin L$, el objetivo es encontrar el vector del retículo $\mathbf{v} \in L$ más cercano a $\mathbf{t}$ \cite{huang2018lattice}.
    
    \item \textbf{Solución de Enteros Cortos (Short Integer Solution - SIS):} Dada una matriz distribuida uniformemente al azar $\mathbf{A} \in \mathbb{Z}_q^{n \times m}$ y un parámetro de cota de norma $\beta > 0$, el problema $\text{SIS}_{n,m,q,\beta}$ consiste en encontrar un vector no nulo $\mathbf{x} \in \mathbb{Z}^m \setminus \{\mathbf{0}\}$ tal que \cite{ma2026laasso, an2023secrethandshakes}:
    \begin{equation}
    \mathbf{A}\mathbf{x} \equiv \mathbf{0} \pmod q \quad \text{y} \quad \|\mathbf{x}\| \le \beta
    \end{equation}
    Resolver SIS equivale a hallar un vector corto en el retículo ortogonal $\Lambda_q^\perp(\mathbf{A})$ \cite{an2023secrethandshakes}.
    
    \item \textbf{Solución de Enteros Cortos en Módulos (Module-SIS - MSIS):} Es la extensión algebraica del problema SIS sobre anillos cociente $\mathcal{R}_q$. Dada una matriz de polinomios $\mathbf{A} \in \mathcal{R}_q^{n \times m}$ y un límite de norma $\beta$, se busca hallar un vector de polinomios $\mathbf{s} \in \mathcal{R}_q^m \setminus \{\mathbf{0}\}$ tal que $\mathbf{A}\mathbf{s} \equiv \mathbf{0} \pmod q$ con $0 < \|\mathbf{s}\| \le \beta$ \cite{ma2026laasso}.
\end{itemize}

\paragraph{Fundamento de Reducción del Peor Caso al Caso Promedio y Resistencia Cuántica}

Los esquemas criptográficos tradicionales como RSA, ECC y ElGamal basan su seguridad en problemas algebraicos clásicos (factorización de enteros y logaritmo discreto) \cite{huang2018lattice, nguyen2019high, rewal2023quantumsafe, mukil2026implementing}. No obstante, en 1994, Peter Shor formuló un algoritmo cuántico capaz de resolver el logaritmo discreto y la factorización entera en tiempo polinomial $O(n^3)$, vulnerando la infraestructura de seguridad pre-cuántica \cite{huang2018lattice, rewal2023quantumsafe, mukil2026implementing}.

Frente a esta amenaza, Ajtai (1996) introdujo la criptografía basada en retículos \cite{huang2018lattice, rewal2023quantumsafe}. Su propiedad más revolucionaria reside en las reducciones de seguridad del peor caso al caso promedio (\textit{worst-case to average-case reductions}) \cite{rewal2023quantumsafe, lai2025leaky, an2023secrethandshakes}: la capacidad de resolver instancias promedio elegidas al azar de problemas como SIS o LWE garantiza la capacidad de resolver instancias en el peor caso posible de problemas geométricos subyacentes como SVP o SIVP sobre cualquier retículo de alta dimensión \cite{rewal2023quantumsafe, an2023secrethandshakes, lai2025leaky}. Dado que no existen algoritmos cuánticos conocidos capaces de resolver las aproximaciones de SVP o CVP en tiempo subexponencial, los retículos ofrecen un fundamento matemático resistente a ataques de la era poscuántica \cite{huang2018lattice, nguyen2019high, rewal2023quantumsafe, ma2026laasso, mukil2026implementing, yang2025multiserver, lombardi2022postquantum}.


\subsubsection{El Problema Learning with Errors (LWE) Estándar}

\paragraph{Formulación Matemática: Versiones de Búsqueda y Decisión}

El problema \textit{Learning with Errors} (LWE), introducido originalmente por Oded Regev \cite{regev2009lwe}, constituye una de las primitivas fundamentales de la criptografía basada en retículos poscuánticos \cite{mukil2026implementing, lai2025leaky}. Matemáticamente, LWE puede entenderse como el problema de resolver un sistema de ecuaciones lineales sobre determinado e ruidoso sobre el cuerpo/anillo finito $\mathbb{Z}_q = \mathbb{Z}/q\mathbb{Z}$, donde $q \ge 2$ representa un módulo entero \cite{regev2009lwe, yang2025multiserver}.

Para definir formalmente el problema, sea $n \in \mathbb{N}$ la dimensión del vector secreto no revelado $\mathbf{s} \in \mathbb{Z}_q^n$, y sea $\chi$ una distribución de probabilidad con soporte sobre los enteros $\mathbb{Z}$ (frecuentemente una distribución Gaussiana discreta $\chi_\sigma$ de media cero y desviación estándar $\sigma$) \cite{nguyen2019high, yang2025multiserver, lai2025leaky}. Una muestra individual de LWE se genera seleccionando un vector de coeficientes uniformemente al azar $\mathbf{a} \leftarrow \mathbb{Z}_q^n$ y muestreando un escalar de ruido $e \leftarrow \chi$, produciendo el par $(\mathbf{a}, b) \in \mathbb{Z}_q^n \times \mathbb{Z}_q$ definido por:

\begin{equation}
b = \mathbf{a}^T \mathbf{s} + e \pmod q
\end{equation}

Cuando se recolectan $m$ muestras independientes en formato matricial, la matriz pública se denota como $\mathbf{A} = [\mathbf{a}_1 | \mathbf{a}_2 | \dots | \mathbf{a}_m] \in \mathbb{Z}_q^{n \times m}$ y el vector de errores aleatorios como $\mathbf{e} = (e_1, e_2, \dots, e_m)^T \leftarrow \chi^m$. El sistema completo de ecuaciones ruidosas se expresa algebraicamente como:

\begin{equation}
\mathbf{b}^T = \mathbf{s}^T \mathbf{A} + \mathbf{e}^T \pmod q
\end{equation}

El problema LWE se formaliza analíticamente en dos variantes computacionales fundamentales \cite{regev2009lwe, lombardi2022postquantum, lai2025leaky}:

\begin{itemize}
    \item \textbf{LWE en versión Búsqueda (Search-LWE):} Dada una matriz de coeficientes aleatorios $\mathbf{A} \in \mathbb{Z}_q^{n \times m}$ y el vector de muestras ruidosas $\mathbf{b}^T = \mathbf{s}^T \mathbf{A} + \mathbf{e}^T \pmod q$, el objetivo computacional consiste en determinar de forma exacta el vector secreto $\mathbf{s} \in \mathbb{Z}_q^n$.
    
    \item \textbf{LWE en versión Decisión (Decision-LWE):} El objetivo de un algoritmo probabilístico de tiempo polinomial (PPT) es distinguir con ventaja no despreciable entre la distribución de pares genuinos LWE $(\mathbf{A}, \mathbf{s}^T \mathbf{A} + \mathbf{e}^T \pmod q)$ y una distribución uniformemente aleatoria $(\mathbf{A}, \mathbf{u}^T) \in \mathbb{Z}_q^{n \times m} \times \mathbb{Z}_q^m$.
\end{itemize}

Intuición técnica: Sin la adición de la variable de error $e \leftarrow \chi$, el sistema lineal $\mathbf{b}^T = \mathbf{s}^T \mathbf{A} \pmod q$ sería un problema algebraico trivial resoluble mediante la eliminación Gaussiana clásica en tiempo polinomial $O(n^3)$. La introducción del ruido aleatorio acotado $e$ destruye la estructura algebraica lineal exacta del sistema, convirtiéndolo en un problema geométricamente intrincado de retículos multidimensionales \cite{regev2009lwe, mukil2026implementing}.

\paragraph{Fundamentación Teórica: Reducción del Peor Caso al Caso Promedio}

El aspecto teórico más relevante de LWE radica en la existencia de una reducción rigurosa que conecta la dureza del problema en el \textit{caso promedio} (average-case) con la intratabilidad de problemas geométricos de retículos en el \textit{peor caso} (worst-case) \cite{regev2009lwe, rewal2023quantumsafe, lai2025leaky}.

Regev \cite{regev2009lwe} demostró formalmente mediante una reducción cuántica que para un módulo primo $q$ y una distribución de error Gaussiana $\chi_\sigma$ de ancho $\alpha = \sigma/q = \omega(\sqrt{\log n}/n)$, si existe un algoritmo eficiente capaz de resolver Decision-LWE en el caso promedio para instancias aleatorias $(\mathbf{A}, \mathbf{b})$, entonces existe un algoritmo cuántico capaz de aproximar problemas geométricos difíciles en el peor caso sobre cualquier retículo de dimensión $n$ (tales como la versión de aproximación polinomial del \textit{Shortest Independent Vectors Problem} $\text{SIVP}_\gamma$ o $\text{SVP}_\gamma$) \cite{regev2009lwe, rewal2023quantumsafe}.

Esta garantía implica que una instancia individual de LWE generada al azar no posee ``casos d\'ebiles'', a diferencia de los esquemas criptográficos tradicionales (RSA o ECC), cuya seguridad depende únicamente de la dificultad promedio de factorizar enteros o calcular logaritmos discretos en grupos específicos \cite{rewal2023quantumsafe, mukil2026implementing, lombardi2022postquantum}.

\paragraph{Parámetros Críticos y Trade-offs de Seguridad y Rendimiento}

El nivel de seguridad de la primitiva LWE y su eficiencia práctica en la construcción de protocolos están gobernados por el balance directo entre cuatro parámetros interdependientes \cite{regev2009lwe, yang2025multiserver, lai2025leaky}:

\begin{itemize}
    \item \textbf{Dimensión Secreta ($n$):} Determina el rango espacial del vector secreto $\mathbf{s} \in \mathbb{Z}_q^n$. Governa la dureza asintótica del problema frente a algoritmos de reducción de bases de retículos (ej. BKZ o algoritmos de cribado \textit{Sieve}) \cite{mukil2026implementing}. Aumentar $n$ incrementa exponencialmente la resistencia cuántica, pero eleva cuadráticamente el tamaño de las matrices públicas $\mathbf{A} \in \mathbb{Z}_q^{n \times m}$ \cite{yang2025multiserver, ma2026laasso}.
    
    \item \textbf{Módulo Aritmético ($q$):} Define la aritmética modular sobre la cual se realizan todas las sumas y multiplicaciones. La relación entre el módulo $q$ y la magnitud del error $\sigma$ (la razón módulo-ruido $q/\sigma$) es crítica. Un ratio $q/\sigma$ elevado simplifica la corrección y reconciliación de errores en el receptor, pero reduce la dureza geométrica de LWE \cite{yang2025multiserver, rewal2023quantumsafe}.
    
    \item \textbf{Número de Muestras ($m$ o $Q$):} Representa el número de ecuaciones lineales expuestas al adversario. Si el número de consultas $m$ es ilimitado o excesivamente grande en comparación con $n$, la seguridad teórica puede degradarse si no se acota la norma de los ruidos \cite{lombardi2022postquantum, lai2025leaky}. En aplicaciones reales, $m$ se parametriza linealmente con la dimensión secreta ($m = \text{poly}(n)$) \cite{yang2025multiserver}.
    
    \item \textbf{Distribución de Error ($\chi$):} Define la distribución aleatoria sobre $\mathbb{Z}$ de los términos $e$. Mientras que el análisis teórico clásico utiliza distribuciones Gaussianas discretas $\chi_\sigma$, implementaciones de alto rendimiento reemplazan la Gaussiana por la \textit{Distribución Binomial Centrada} (CBD con parámetro $\eta$), la cual permite realizar el muestreo en tiempo constante $O(1)$\cite{yang2025multiserver, mukil2026implementing}.
\end{itemize}



\subsubsection{Variantes Algebraicas y Estructurales del Problema LWE}

\paragraph{Ring-LWE (RLWE) y la Aceleración mediante la Transformada de Teoría de Números (NTT)}
El problema \textit{Learning with Errors} (LWE) estándar impone requerimientos de almacenamiento y sobrecarga computacional prohibitivos para muchos sistemas prácticos, dado que su matriz pública exige un espacio en memoria de orden $O(n^2)$. Para resolver esta limitación, se introdujo \textit{Ring-Learning with Errors} (RLWE), una variante algebraicamente estructurada que traslada la aritmética desde el cuerpo finito $\mathbb{Z}_q$ hacia un anillo polinomial cociente $\mathcal{R}_q = \mathbb{Z}_q[x]/\langle f(x) \rangle$ \cite{nguyen2019high, mukil2026implementing}. En esta estructura, $f(x) = x^n + 1$ se fija habitualmente como un polinomio ciclotómico irreducible de grado $n$ (donde $n$ es una potencia de 2) y $q$ es un módulo primo. 

La matriz pública $\mathbf{A} \in \mathbb{Z}_q^{m \times n}$ de LWE se reemplaza en RLWE por un único polinomio $a(x) \leftarrow \mathcal{R}_q$ elegido uniformemente al azar, reduciendo drásticamente el espacio de la clave pública de $O(n^2)$ a $O(n)$ \cite{nguyen2019high}. El secreto $s(x)$ y el error $e(x)$ se definen como polinomios en $\mathcal{R}_q$, muestreando los coeficientes del error desde una distribución discreta Gaussiana $\chi_\sigma$ (o una distribución binomial centrada equivalente en implementaciones prácticas). Una muestra del problema decisional RLWE se genera algebraicamente como:

\begin{equation}
b(x) = a(x) \cdot s(x) + e(x) \pmod q
\end{equation}

La ventaja computacional determinante de RLWE radica en la optimización de la multiplicación polinomial mediante la Transformada de Teoría de Números (NTT), la cual reduce la complejidad temporal de la convolución discreta de $O(n^2)$ a $O(n \log n)$ \cite{nguyen2019high}. Para evitar la latencia y la sobrecarga de memoria asociadas al rellenado de ceros (\textit{zero-padding}), se utiliza la \textit{convolución envuelta negativa}. Si el anillo admite una raíz primitiva $n$-ésima de la unidad $\omega \in \mathbb{Z}_q$, se define el factor de escala $\psi \equiv \sqrt{\omega} \pmod q$. Los coeficientes de los polinomios $a(x)$ y $b(x)$ se preprocesan escalando su $i$-ésimo término por $\psi^i$ (obteniendo los vectores de coeficientes modificado $a'$ y $b'$). De este modo, la multiplicación en $\mathcal{R}_q$ se reduce a un producto elemento a elemento ($\odot$) en el dominio transformado:

\begin{equation}
c' = \text{INTT}_{\omega}^n \left( \text{NTT}_{\omega}^n(a') \odot \text{NTT}_{\omega}^n(b') \right)
\end{equation}

Este enfoque minimiza la latencia computacional y el área de silicio requerida en hardware, convirtiendo a RLWE en un pilar fundamental de la criptografía poscuántica \cite{nguyen2019high, mukil2026implementing}.

\paragraph{Module-LWE (MLWE)}
Aunque RLWE logra una reducción dramática en la sobrecarga espacial y temporal, la estructura de retículo ideal subyacente impone simetrías algebraicas rigurosas que suscitan ciertas prevenciones teóricas en materia de seguridad. \textit{Module-Learning with Errors} (MLWE) surge como un marco algebraico intermedio que logra un balance óptimo entre la seguridad estricta de LWE sobre retículos no estructurados y la alta eficiencia de RLWE \cite{mukil2026implementing, yang2025multiserver}. En lugar de trabajar con escalares o con un único polinomio de grado elevado, MLWE emplea matrices de dimensión reducida $k \times k$ (donde $k \ge 1$ es la dimensión del módulo) cuyos elementos son polinomios pertenecientes al anillo $\mathcal{R}_q$.

Bajo este paradigma, la matriz pública se define como $\mathbf{A} \in \mathcal{R}_q^{k \times k}$, mientras que el vector secreto $\mathbf{s} \in \mathcal{R}_q^k$ y el vector de error $\mathbf{e} \in \mathcal{R}_q^k$ están formados por $k$ polinomios independientes muestreados de la distribución de ruido. Una muestra del problema decisional de módulo (DMLWE) se expresa como la tupla:

\begin{equation}
(\mathbf{A}, \, \mathbf{b} = \mathbf{A}\mathbf{s} + \mathbf{e} \pmod q) \in \mathcal{R}_q^{k \times k} \times \mathcal{R}_q^k
\end{equation}

Además de atenuar las simetrías de los ideales de anillo, MLWE aporta la propiedad operativa de la escalabilidad modular \cite{yang2025multiserver}: permite incrementar el nivel de seguridad del sistema (por ejemplo, ajustándose a las categorías 1, 3 o 5 de estandarización del NIST) modificando únicamente el parámetro matricial $k$ ($k=2, 3, 4$), sin necesidad de modificar el polinomio base $f(x) = x^{256} + 1$ ni rediseñar los coprocesadores hardware concebidos para la NTT. Esta flexibilidad ha llevado a que estándares internacionales como \textit{CRYSTALS-Kyber} fundamenten la seguridad de su mecanismo de encapsulación de claves (KEM) directamente en la dureza de MLWE \cite{mukil2026implementing}.

\paragraph{Leaky LWE (LWE con Fugas)}
El modelo \textit{Leaky LWE} formaliza la seguridad del problema en escenarios de canal lateral o ejecución adversaria avanzada, asumiendo que el atacante puede obtener información residual (fugas) sobre el estado interno del secreto y del error de forma semi-adaptativa \cite{lai2025leaky}. En este esquema, además de la tupla pública $(\mathbf{A}, \mathbf{b})$, el adversario recibe una medida de fuga ruidosa expresada como:

\begin{equation}
\mathbf{l}^T = (\mathbf{s}^T, \mathbf{e}^T)\mathbf{L} + \mathbf{f}^T
\end{equation}

donde $\mathbf{L}$ representa una matriz de fuga de norma acotada (que puede ser seleccionada semi-adaptativamente por el atacante tras observar la matriz pública $\mathbf{A}$) y $\mathbf{f}$ es un vector de error independiente asociado al canal de fuga \cite{lai2025leaky}. 

Históricamente, garantizar la seguridad en presencia de fugas requería recurrir a la técnica de ``inundación de ruido estadístico'' (\textit{noise flooding}), la cual obligaba a emplear módulos $q$ de tamaño superpolinomial ($q > 2^{128}$) para ahogar la información filtrada, penalizando fuertemente el rendimiento del sistema. \textit{Leaky LWE} demuestra formalmente una reducción de seguridad estrecha desde LWE estándar, probando que la presencia del vector $\mathbf{l}$ no trivializa la extracción del secreto $\mathbf{s}$ ni compromete la pseudoaleatoriedad de la muestra \cite{lai2025leaky}. Esto permite a los diseñadores utilizar módulos $q$ polilogarítmicos mucho más pequeños (donde $q$ escala de manera proporcional a $\sqrt{Q}$, siendo $Q$ el número de consultas de fuga del atacante, permitiendo módulos prácticos del orden de $q \approx 2^{32}$). Este avance resulta determinante para eliminar la sobrecarga computacional en protocolos complejos como las pruebas de conocimiento cero (ZKP), las firmas umbral y los esquemas criptográficos basados en identidad \cite{lai2025leaky}.

\paragraph{Binary-Secret LWE (BLWE)}
\textit{Binary-Secret Learning with Errors} (BLWE) constituye una especialización de LWE en la que la selección de las componentes del vector secreto $\mathbf{s}$ se restringe estrictamente al dominio binario $\{0, 1\}^n$ \cite{an2023secrethandshakes}. Restringir el secreto al alfabeto binario sustituye las multiplicaciones vectoriales de alto costo por sumas vectoriales condicionales simples y reduce sensiblemente la huella de memoria de las claves privadas, resultando altamente atractivo en la construcción de protocolos para dispositivos de recursos muy restringidos y en esquemas de reconocimiento mutuo anónimo (\textit{secret handshakes}) \cite{an2023secrethandshakes}.

\paragraph{Learning with Rounding (LWR) y Ring-LWR}
A diferencia de LWE o RLWE, cuyas pruebas de reducción de seguridad dependen de la adición estocástica de un término de error $e$ muestreado en tiempo de ejecución desde una distribución discreta Gaussiana, la variante \textit{Learning with Rounding} (LWR) elimina la necesidad de generar ruido aleatorio y lo sustituye por una operación de redondeo modular determinista \cite{mukil2026implementing}. Específicamente, el error de cuantización inducido al escalar la muestra desde un módulo mayor $q$ hacia un módulo menor $p < q$ mediante la función determinista $\lfloor \frac{p}{q} \cdot x \rfloor$, reemplaza funcionalmente al ruido gaussiano. 

Esquemas derivados como \textit{Lizard} y \textit{Ring-Lizard} combinan los beneficios de espacio de la estructura de anillo (Ring-LWE) con el redondeo determinista de LWR \cite{mukil2026implementing}. Al suprimir por completo los generadores de entropía y el muestreo de gaussianas discretas en tiempo de ejecución, estas variantes reducen radicalmente el consumo energético, resultando óptimas para arquitecturas del Internet de las Cosas (IoT) y entornos con restricciones severas de cómputo \cite{mukil2026implementing}.

\paragraph{Pairing with Error (PWE) y Decision PWE (DPWE)}
Las suposiciones de dureza \textit{Pairing with Error} (PWE) y su variante decisional (DPWE) se diseñan e introducen primordialmente para respaldar las demostraciones formales de seguridad en protocolos de Intercambio de Claves Autenticado por Contraseña (PAKE) en entornos poscuánticos \cite{yang2025multiserver, rewal2023quantumsafe}. En estos esquemas, las partes intervinientes calculan productos LWE ruidosos y utilizan mecanismos avanzados de reconciliación modular de errores (como los esquemas de Peikert o Ding) para corregir las discrepancias producidas por el error $e$ y concertar una clave simétrica idéntica. 

La hipótesis PWE/DPWE establece formalmente que la clave derivada mediante dicha función de reconciliación permanece computacionalmente indistinguible de un valor uniforme aleatorio para cualquier observador externo que intercepte tanto los mensajes LWE públicos como la señal de reconciliación enviada por el canal abierto \cite{yang2025multiserver, rewal2023quantumsafe}.

\subsubsection{Muestreo y Gestión del Ruido: De la Distribución Gaussiana a la Distribución Binomial Centrada (CBD)}

\paragraph{Distribución Gaussiana Discreta y Algoritmos Clásicos de Muestreo}

En la criptografía basada en retículos, la seguridad del problema \textit{Learning with Errors} (LWE) y sus variantes algebraicas depende intrínsecamente de la inyección de un término de error no determinista e impredecible \cite{regev2009lwe, yang2025multiserver}. Tradicionalmente, este ruido se modela mediante una distribución Gaussiana discreta $\chi_\sigma$ sobre el anillo de enteros $\mathbb{Z}$ o sobre un retículo $n$-dimensional $\Lambda \subset \mathbb{R}^n$, centrada en el origen y parametrizada por la desviación estándar $\sigma > 0$ (o parámetro Gaussiano $s = \sigma \sqrt{2\pi}$) \cite{nguyen2019high, huang2018lattice, rewal2023quantumsafe}.

Para un centro $\mathbf{c} \in \mathbb{R}^n$ y un parámetro $s > 0$, la función de densidad de probabilidad Gaussiana discreta asigna a cada vector $\mathbf{x} \in \Lambda$ una probabilidad proporcional a su peso exponencial euclídeo \cite{huang2018lattice, lai2025leaky}:

\begin{equation}
D_{\Lambda, s, \mathbf{c}}(\mathbf{x}) = \frac{\exp\left(-\pi \frac{\|\mathbf{x} - \mathbf{c}\|^2}{s^2}\right)}{\sum_{\mathbf{y} \in \Lambda} \exp\left(-\pi \frac{\|\mathbf{y} - \mathbf{c}\|^2}{s^2}\right)}
\end{equation}

En la literatura, la generación de muestras pertenecientes a la distribución Gaussiana discreta se realiza históricamente mediante dos familias de algoritmos principales \cite{huang2018lattice, nguyen2019high}:

\begin{itemize}
    \item \textbf{Algoritmo de Knuth-Yao:} Diseñado para el muestreo de ruido escalar sobre enteros, representa la distribución teórica a través de un Árbol de Grafos de Distribución Discreta (\textit{Discrete Distribution Graph} - DDG) \cite{nguyen2019high}. El algoritmo recorre la estructura de árbol consumiendo bits aleatorios de entropía para seleccionar el coeficiente final. Aunque Knuth-Yao demuestra una eficiencia cercana a la entropía teórica ideal (optimizando el número de bits aleatorios requeridos), depende de tablas de búsqueda precalculadas y de un número variable de comparaciones en función de los bits muestreados \cite{nguyen2019high}.
    
    \item \textbf{Algoritmo de Muestreo de Trapdoor GPV (\texttt{SampleD}):} Desarrollado por Gentry, Peikert y Vaikuntanathan, se utiliza para extraer vectores cortos en retículos de la forma $\Lambda_q^\perp(\mathbf{A})$ mediante el uso de una base de trampa corta $\mathbf{T}_A$ \cite{huang2018lattice}. Dado un vector objetivo $\mathbf{t} \in \mathbb{Z}_q^n$ y una matriz pública $\mathbf{A} \in \mathbb{Z}_q^{n \times m}$, el algoritmo $\texttt{SampleD}(\mathbf{T}_A, s, \mathbf{c}, \mathbf{t})$ genera un vector aleatorio $\mathbf{x} \in \mathbb{Z}^m$ distribuido según $D_{\Lambda, s, \mathbf{c}}$ que satisface la relación $\mathbf{A}\mathbf{x} \equiv \mathbf{t} \pmod q$ \cite{huang2018lattice}. La unidireccionalidad (\textit{one-wayness}) de este algoritmo fundamenta la seguridad de esquemas de autenticación basados en trampas \cite{huang2018lattice}.
\end{itemize}

A pesar de sus sólidas propiedades estadísticas, el muestreo Gaussiano discreto clásico presenta severas limitaciones en entornos físicos de ejecución \cite{mukil2026implementing, yang2025multiserver}. La necesidad de evaluar funciones trascendentes, usar tablas de búsqueda dependientes de los datos o realizar iteraciones variables en algoritmos como Knuth-Yao induce variaciones en los tiempos de respuesta y en los perfiles de consumo energético, exponiendo al dispositivo a ataques de canal lateral (\textit{Side-Channel Attacks} - SCA) y ataques de temporización (\textit{timing attacks}) \cite{mukil2026implementing}.

\paragraph{Muestreador de Distribución Binomial Centrada (CBD)}

Con el objetivo de eliminar la sobrecarga computacional del muestreo Gaussiano y garantizar un comportamiento algebraico homogéneo, las construcciones modernas de retículos (como Kyber y Module-LWE) adoptan el Muestreador de Distribución Binomial Centrada (\textit{Centered Binomial Distribution} - CBD) \cite{yang2025multiserver, mukil2026implementing}.

La distribución CBD, denotada como $\psi_\eta$, aproxima la curva de una distribución Gaussiana discreta centrada en cero con parámetro de cota entero $\eta \in \mathbb{Z}^+$ \cite{yang2025multiserver}. Para generar un coeficiente de error $e \in [-\eta, \eta]$, el algoritmo toma $2\eta$ bits aleatorios independientes e idénticamente distribuidos $a_0, a_1, \dots, a_{\eta-1}, b_0, b_1, \dots, b_{\eta-1} \in \{0, 1\}$ y calcula la diferencia de sus pesos de Hamming \cite{yang2025multiserver}:

\begin{equation}
e = \sum_{i=0}^{\eta-1} a_i - \sum_{i=0}^{\eta-1} b_i
\end{equation}

Por simetría, la esperanza matemática de la distribución es $E[e] = 0$, y su varianza viene dada por $\sigma^2 = \eta / 2$, correspondiente a una desviación estándar de $\sigma = \sqrt{\eta/2}$ \cite{yang2025multiserver}.

En la parametrización estándar de Module-LWE, se selecciona el parámetro $\eta = 2$, lo que fija la desviación estándar en $\sigma = 1$ y restringe el rango de los coeficientes del error al conjunto discreto corto $\{-2, -1, 0, 1, 2\}$ \cite{yang2025multiserver}. El muestreo de un polinomio de error de grado $d=256$ requiere únicamente la extracción de $2 \cdot \eta \cdot d = 1024$ bits de aleatoriedad uniforme, procesados mediante operaciones de bits paralelas sobre palabras del microprocesador \cite{yang2025multiserver}.

\paragraph{Fundamento de Inmunidad Física y Aplicabilidad al Prototipo}

La sustitución de la distribución Gaussiana discreta por el algoritmo CBD con $\eta = 2$ representa una decisión de diseño fundamental para garantizar la factibilidad física y la seguridad de un prototipo de autenticación \cite{yang2025multiserver, mukil2026implementing}:

\begin{itemize}
    \item \textbf{Ejecución en Tiempo Constante Estricto $O(1)$:} A diferencia del algoritmo Knuth-Yao o del muestreo por rechazo Gaussiano, la CBD no contiene bucles condicionales de duración indeterminada ni saltos de ejecución (\textit{if-else}) dependientes de datos secretos \cite{yang2025multiserver, mukil2026implementing}. Cada coeficiente se calcula mediante una secuencia fija y acotada de instrucciones booleanas a nivel de procesador (\texttt{AND}, \texttt{XOR}, desplazamientos de bits y conteo de peso de Hamming), asegurando que el tiempo de cómputo sea idéntico e independiente de los valores muestreados \cite{yang2025multiserver}.
    
    \item \textbf{Eliminación de Tablas de Búsqueda y Protección contra SCA:} Los accesos a memoria indexados por secretos en tablas de probabilidad Gaussiana constituyen una fuente principal de fugas por emisiones electromagnéticas y análisis diferencial de potencia (DPA) \cite{mukil2026implementing}. Al operar directamente sobre registros del procesador mediante operaciones binarias, la CBD genera una firma de consumo de energía plana y predecible, blindando el coprocesador de la tarjeta inteligente contra ataques de canal lateral de traza única (\textit{single-trace side-channel attacks}) \cite{mukil2026implementing}.
    
    \item \textbf{Eficiencia y Compactación de Memoria:} El algoritmo CBD elimina por completo la necesidad de almacenar tablas de precisión en coma flotante en la memoria sin afectar la indistinguibilidad estadística requerida para sostener la prueba de seguridad del protocolo \cite{yang2025multiserver, lai2025leaky}.
\end{itemize}
\subsubsection{Mecanismos de Reconciliación de Errores y Mapeo Criptográfico}

\paragraph{El Problema de la Dispersión de Ruido en la Extracción de Claves Deterministas}

En la criptografía basada en el problema \textit{Learning with Errors} (LWE) y sus variantes algebraicas (RLWE y MLWE), las operaciones interactivas entre dos entidades (cliente y servidor) involucran la inyección deliberada de errores Gaussianos o binomiales cortos para garantizar la seguridad frente a adversarios clásicos y cuánticos \cite{regev2009lwe, yang2025multiserver}. Sin embargo, esta perturbación estocástica introduce un desafío algebraico fundamental: cuando ambas partes calculan un producto escalar o polinomial compartido sobre el anillo $\mathcal{R}_q = \mathbb{Z}_q[x]/\langle x^d + 1 \rangle$, los valores resultantes $w \in \mathbb{Z}_q$ y $w' \in \mathbb{Z}_q$ no son exactamente iguales, sino que difieren por un término de error acumulado no nulo $e \in \mathbb{Z}_q$, tal que $w' = w + e \pmod q$ \cite{yang2025multiserver, rewal2023quantumsafe}.

En un esquema de autenticación o acuerdo de claves determinista, esta pequeña discrepancia aritmética impediría que ambas entidades obtengan exactamente la misma secuencia de bits de secreto, provocando fallos continuos de autenticación \cite{yang2025multiserver, mukil2026implementing}. Para resolver esta dispersión sin transmitir información en texto claro que comprometa la clave, se emplean técnicas de mapeo discreto y mecanismos de reconciliación de errores \cite{yang2025multiserver, rewal2023quantumsafe}.

\paragraph{Codificación de Bits en Coeficientes Modulares}

Para procesar datos binarios discretos dentro del espacio modular $\mathbb{Z}_q$, se requiere una función de codificación que mapee los símbolos del alfabeto binario $\{0, 1\}$ hacia coeficientes polinomiales con la máxima separación modular posible \cite{nguyen2019high, mukil2026implementing}.

Dado un módulo entero $q$, la función de codificación mapea cada bit binario $b_i \in \{0, 1\}$ a un coeficiente polinomial $m_i \in \mathbb{Z}_q$ mediante la siguiente regla determinista \cite{nguyen2019high}:

\begin{equation}
m_i = \text{Encode}(b_i) = \begin{cases} 0 & \text{si } b_i = 0 \\ \frac{q-1}{2} & \text{si } b_i = 1 \end{cases} \pmod q
\end{equation}

Esta asignación ubica la representación del bit $1$ en el punto medio del espectro modular $\mathbb{Z}_q$ (a distancia máxima de $0$), lo que maximiza la tolerancia geométrica del sistema frente a las desviaciones causadas por el ruido LWE durante el proceso de descifrado y reconciliación \cite{nguyen2019high, rewal2023quantumsafe}.

\paragraph{Mecanismo de Reconciliación de Peikert y Funciones de Redondeo}

Para permitir que dos entidades concilien el mismo bit secreto a partir de valores ruidosos aproximados $w, w' \in \mathbb{Z}_q$, se utiliza el mecanismo de reconciliación de errores de Peikert \cite{yang2025multiserver}. Supóngase que $q$ es un entero modular. Se definen los siguientes subconjuntos de partición sobre $\mathbb{Z}_q$ \cite{yang2025multiserver}:

\begin{equation}
I_0 = \left\{0, 1, \dots, \left\lfloor \frac{q}{4} \right\rfloor - 1\right\}, \quad I_1 = \left\{-\left\lfloor \frac{q}{4} \right\rfloor, \dots, -1\right\}, \quad \sigma = \left[-\left\lfloor \frac{q}{8} \right\rfloor, \left\lfloor \frac{q}{8} \right\rfloor\right) \cap \mathbb{Z}_q
\end{equation}

A partir de esta partición, Peikert define tres funciones matemáticas fundamentales \cite{yang2025multiserver}:

\begin{enumerate}
    \item \textbf{Función de Redondeo Cruzado (\textit{Cross-Rounding}):} Extrae un bit de señal de sincronización $v \in \{0, 1\}$ determinando en qué región del cuadrante modular se encuentra el valor:
    \begin{equation}
    \langle x \rangle_{q, 2} = \left\lfloor \frac{4}{q} \cdot x \right\rfloor \pmod 2
    \end{equation}
    
    \item \textbf{Función de Redondeo Modular:} Extrae el bit de clave altamente significativo redondeando el valor modular:
    \begin{equation}
    \lfloor x \rfloor_{q, 2} = \left\lfloor \frac{2}{q} \cdot x \right\rfloor \pmod 2
    \end{equation}
    
    \item \textbf{Función de Reconciliación Determinista $rec(w, v)$:} Utiliza el valor ruidoso $w$ y la señal de ayuda $v = \langle w' \rangle_{q, 2}$ enviada por la otra parte para recuperar de forma exacta el bit redondeado original:
    \begin{equation}
    rec(w, v) = \begin{cases} 0 & \text{si } w \in I_v + \sigma \pmod q \\ 1 & \text{en otro caso} \end{cases}
    \end{equation}
\end{enumerate}

Garantía de Corrección: Si dos valores modulares difieren por un error acotado $w' = w + e \pmod q$ con $|e| < q/8$, entonces el mecanismo de Peikert garantiza matemáticamente que $\lfloor w \rfloor_{q, 2} = rec(w', \langle w \rangle_{q, 2})$, logrando una concordancia de bits del 100\% entre ambas partes \cite{yang2025multiserver}. De forma complementaria, variantes como la función auxiliar modular $f_2(u, d) = \left(u + d \cdot \frac{q-1}{2}\right) \bmod q \bmod 2$ permiten obtener resultados idénticos ajustando la paridad modular \cite{rewal2023quantumsafe}.
\subsubsection{Tolerancia a la Fuga de Información y Análisis de Resistencia: Leaky LWE y Reducción del Módulo}

\paragraph{La Problemática de la Inundación Estadística de Ruido (\textit{Noise Flooding})}

En la teoría de protocolos criptográficos y pruebas de conocimiento cero (\textit{Zero-Knowledge Proofs} - ZKP) basados en la dureza de problemas de retículos, la demostración formal de la propiedad de conocimiento cero o de resistencia ante consultas adaptativas ha dependido históricamente de la técnica conocida como \textit{inundación estadística de ruido} (\textit{noise flooding}, \textit{noise drowning} o \textit{noise smudging}) \cite{goldwasser1989knowledge, lai2025leaky}.

Cuando un probador emite una respuesta durante un protocolo interactivo o no interactivo, dicha respuesta combina linealmente los elementos secretos con un ruido de enmascaramiento. Algebraicamente, la respuesta actúa como una fuga lineal ruidosa de la forma $\mathbf{l}^T = (\mathbf{s}^T, \mathbf{e}^T)\mathbf{L} + \mathbf{f}^T$, donde $\mathbf{s} \in \mathbb{Z}_q^n$ representa la clave secreta, $\mathbf{e} \in \mathbb{Z}_q^m$ es el vector de error de LWE, $\mathbf{L} \in \mathbb{Z}_q^{(n+m) \times Q}$ es la matriz de desafío del verificador y $\mathbf{f} \in \mathbb{Z}_q^Q$ es el ruido de enmascaramiento \cite{lai2025leaky}. Para asegurar que la distribución de $\mathbf{l}$ sea estadísticamente indistinguible de un valor independiente del secreto, el principio de \textit{noise flooding} exige que el parámetro de dispersión del ruido $\mathbf{f}$ sea superpolinomialmente mayor que la norma del secreto enmascarado \cite{lai2025leaky}:

\begin{equation}
\sigma_f \ge \text{superpoly}(\lambda) \cdot \|(\mathbf{s}^T, \mathbf{e}^T)\mathbf{L}\|
\end{equation}

Esta inundación masiva de ruido impone un grave inconveniente estructural: para evitar desbordamientos aritméticos que causen fallos de decodificación o invaliden las relaciones modulares, el módulo $q$ del sistema debe fijarse en un entero superpolinomialmente gigante, típicamente $q > 2^{128}$ \cite{lai2025leaky}. Para mantener la dureza del problema LWE frente a algoritmos de reducción de retículos ante un módulo $q$ tan elevado, es obligatorio escalar la dimensión del espacio $n$ a valores masivos \cite{regev2009lwe, lai2025leaky}. Como resultado, las claves públicas y las pruebas criptográficas alcanzan tamaños de varios megabytes, limitando su viabilidad en entornos de cómputo restringidos \cite{mukil2026implementing, lai2025leaky}.

\paragraph{Fundamentos de Leaky LWE y Fugas Semi-Adaptativas}

Con el fin de prescindir de la inundación estadística de ruido y sustentar la seguridad en una reducción estrictamente computacional, Lai, Swarnakar y Woo \cite{lai2025leaky} formalizaron la suposición \textit{Leaky LWE} (LWE con fugas). El problema Leaky LWE evalúa la dificultad de distinguir una muestra LWE genuina de una muestra uniformemente aleatoria en presencia de fugas lineales ruidosas introducidas de forma semi-adaptativa \cite{lai2025leaky}.

Formalmente, dado el módulo $q$, las dimensiones $n, m, Q \in \mathbb{N}$ y las distribuciones de error $D_s, D_e, D_f$, el experimento de decisión Leaky LWE se define como sigue \cite{lai2025leaky}:

\begin{enumerate}
    \item Se muestrea una matriz pública uniforme $\mathbf{A} \leftarrow \mathbb{Z}_q^{n \times m}$.
    \item El adversario observa la matriz $\mathbf{A}$ y selecciona de manera \textit{semi-adaptativa} una matriz de fuga de baja norma $\mathbf{L} = \begin{pmatrix} \mathbf{L}_s \\ \mathbf{L}_e \end{pmatrix} \in \mathbb{Z}^{(n+m) \times Q}$ sujeta a una cota de norma espectral \cite{lai2025leaky}.
    \item Se muestrean el secreto $\mathbf{s} \leftarrow D_s$, el error LWE $\mathbf{e} \leftarrow D_e$ y el ruido de fuga $\mathbf{f} \leftarrow D_f$.
    \item El adversario recibe la tupla $(\mathbf{A}, \mathbf{b}^T, \mathbf{l}^T)$, donde $\mathbf{l}^T = (\mathbf{s}^T, \mathbf{e}^T)\mathbf{L} + \mathbf{f}^T \pmod q$, y el vector $\mathbf{b}^T$ es una muestra LWE real $\mathbf{b}^T = \mathbf{s}^T\mathbf{A} + \mathbf{e}^T \pmod q$ (si $b=0$) o un vector aleatorio $\mathbf{b}^T \leftarrow \mathbb{Z}_q^m$ (si $b=1$) \cite{lai2025leaky}.
\end{enumerate}

Lai et al. \cite{lai2025leaky} demostraron que Leaky LWE no es computacionalmente más fácil de resolver que el problema LWE estándar bajo las mismas dimensiones. Mediante transformaciones lineales sobre distribuciones Gaussianas discretas condicionales, los autores construyeron una reducción de seguridad estricta (\textit{tight reduction}) desde LWE clásico hacia Leaky LWE, demostrando que la presencia de fugas ruidosas no reduce la seguridad del problema original \cite{lai2025leaky}.

\paragraph{Cota de Norma de Fuga y Escalamiento del Módulo $q \propto \sqrt{Q}$}

Para que la reducción de seguridad de Leaky LWE sea válida, la matriz de fuga $\mathbf{L}$ elegida semi-adaptativamente debe satisfacer una cota de norma espectral $\beta$ \cite{lai2025leaky}:

\begin{equation}
\|\mathbf{L}^T \mathbf{L}\| \le \beta
\end{equation}

El resultado de esta cota es que la magnitud del ruido no requiere un crecimiento superpolinomial, sino que escala en función de la raíz cuadrada del número total de columnas de fuga o consultas del adversario $Q$ \cite{lai2025leaky}:

\begin{equation}
q \ge c \cdot \sqrt{Q}
\end{equation}

donde $c > 0$ es una constante dependiente de los parámetros del retículo \cite{lai2025leaky}. En aplicaciones donde el número de consultas o transacciones está acotado polinomialmente (por ejemplo, $Q \approx 2^{60}$ a $2^{64}$), el módulo $q$ requiere únicamente un incremento moderado de aproximadamente $30$ bits en lugar de los $128$ bits exigidos por \textit{noise flooding} \cite{lai2025leaky}. Esto permite el uso de módulos compactos de 16 a 32 bits (tales como $q = 3329$) garantizando la reducción formal de seguridad poscuántica \cite{yang2025multiserver, lai2025leaky}.

\paragraph{Análisis de Seguridad Geométrica y Resistencia ante Ataques de Reducción de Retículos (BKZ)}

La reducción del módulo a valores pequeños ($q \in [2^{16}, 2^{32}]$) requiere analizar la seguridad del sistema frente a algoritmos de reducción de bases de retículos, siendo el algoritmo BKZ (\textit{Block-Korkine-Zolotarev}) el estándar de referencia para medir la resistencia geométrica clásica y cuántica \cite{badiezadegan2026complexity, mukil2026implementing}.

Un sistema sobre retículos resultaría vulnerable a ataques BKZ si se fijara simultáneamente un módulo $q$ pequeño y una dimensión espacial reducida \cite{badiezadegan2026complexity}. BKZ opera reduciendo las bases del retículo mediante el desplazamiento de una ventana o bloque de dimensión reducida $\beta_{\text{BKZ}} \le N$ a lo largo del espacio vectorial de dimensión $N$, resolviendo el problema del vector más corto (SVP) exacto dentro de cada sub-bloque mediante algoritmos de cribado cuántico (\textit{quantum sieving}) \cite{badiezadegan2026complexity, mukil2026implementing}. La complejidad temporal de resolver SVP sobre cada bloque de dimensión $\beta_{\text{BKZ}}$ crece de forma exponencial \cite{badiezadegan2026complexity, mukil2026implementing}:

\begin{equation}
T_{\text{BKZ}}(\beta_{\text{BKZ}}) \approx \mathcal{O}\left(2^{c \cdot \beta_{\text{BKZ}}}\right) \quad \text{donde } c \approx 0.265 \text{ a } 0.292
\end{equation}

Para que un atacante no logre extraer el secreto o el error utilizando BKZ sobre un retículo de dimensión $N$, debe ocurrir que al alcanzar un factor de calidad geométrico $\delta_0$ suficientemente pequeño se le exija fijar un tamaño de bloque $\beta_{\text{BKZ}}$ elevado \cite{badiezadegan2026complexity}. Un bloque $\beta_{\text{BKZ}} \ge 400$ demanda una cantidad de operaciones cuánticas del orden de $T_{\text{BKZ}} \ge 2^{160}$, valor que supera holgadamente el umbral de 128 bits de seguridad poscuántica (NIST Nivel 1) y garantiza que el sistema permanezca inexpugnable ante ataques geométricos de reducción de bases \cite{badiezadegan2026complexity, mukil2026implementing, yang2025multiserver}.

\paragraph{Comportamiento de la Matriz de Desafío $\mathbf{L}$ y Rechazo de Consultas Maliciosas}

El modelo de Leaky LWE requiere analizar el comportamiento de la matriz de desafío o fuga $\mathbf{L}$ en escenarios donde un verificador o adversario active consultas en el protocolo:

\begin{itemize}
    \item \textbf{Impacto de una Matriz $\mathbf{L}$ de Norma Elevada:} Si un adversario intentara introducir una matriz de fuga $\mathbf{L}$ con coeficientes grandes (alta norma espectral), el producto $\mathbf{s}\mathbf{L}$ amplificaría geométricamente los valores de la clave secreta $\mathbf{s}$ por encima de la magnitud del ruido de enmascaramiento $\mathbf{f}$ \cite{lai2025leaky}. Como establecieron Alperin-Sheriff y Peikert, si la norma de $\mathbf{L}$ no estuviera acotada, el adversario podría calcular la pseudoinversa de $\mathbf{L}$ sobre los enteros y recuperar algebraicamente el secreto $\mathbf{s}$ \cite{lai2025leaky}.
    
    \item \textbf{Mecanismo de Rechazo y Aborto de Transacción:} En los protocolos criptográficos, la matriz de desafío $\mathbf{L}$ o la respuesta calculada están sujetas a una verificación estricta de norma espectral y norma euclidiana ($\|\mathbf{z}\| \le \beta$) \cite{lombardi2022postquantum, lai2025leaky}. Si el adversario envía un desafío de norma elevada para intentar medir el secreto, el algoritmo de verificación o el muestreo de rechazo del probador detecta que la norma excede el límite $\beta$ y aborta la transacción inmediatamente, evitando responder o anulando la respuesta antes de que se emita información expuesta \cite{lombardi2022postquantum, lai2025leaky}.
    
    \item \textbf{Garantía de Cero Fuga para Norma Baja:} Cuando la matriz de desafío respeta la cota de norma espectral ($\|\mathbf{L}^T \mathbf{L}\| \le \beta$), el Teorema de Leaky LWE garantiza matemáticamente que la respuesta obtenida es computacionalmente indistinguible de ruido aleatorio uniforme \cite{lai2025leaky}. De este modo, la interacción no filtra ningún bit de información utilizable sobre la clave secreta $\mathbf{s}$, preservando la resistencia poscuántica y la propiedad de conocimiento cero del sistema \cite{lombardi2022postquantum, lai2025leaky}.
\end{itemize}
\subsubsection{Protección del Canal de Comunicación: Token Hiding con Kyber y AES}

\paragraph{El Problema del Verificador No Designado y la Intercepción por Terminales Comprometidos}

En los protocolos de autenticación distribuidos y sistemas de inicio de sesión único (\textit{Single Sign-On} - SSO), las credenciales y pruebas criptográficas generadas por el usuario deben transmitirse frecuentemente a través de intermediarios no confiables o redes públicas antes de alcanzar la entidad verificadora final \cite{ma2026laasso}.

Si las pruebas de conocimiento cero (ZKP) o los tokens de autenticación se transmitieran en texto claro a través del canal de comunicación, surgiría el problema del \textit{verificador no designado} \cite{ma2026laasso}. Un cajero automático comprometido mediante soporte malicioso o dispositivos de captura ilegítima (\textit{skimmers}) podría inspeccionar el contenido del token, extraer metadatos de la transacción o rastrear las pruebas enviadas en sesiones sucesivas \cite{ma2026laasso, lombardi2022postquantum}. Esta exposición pública destruye la propiedad de no vinculabilidad (\textit{unlinkability}) y permite a adversarios intermedios perfilarlos hábitos de uso y desplazamientos del cliente dentro de la red bancaria \cite{ma2026laasso, lombardi2022postquantum}.

\paragraph{El Mecanismo de Ocultación de Tokens (\textit{Token Hiding}) mediante Cifrado Híbrido}

Para neutralizar la amenaza de intercepción en nodos intermedios y garantizar que únicamente el servidor central designado pueda validar la autenticación, el protocolo adopta la técnica de \textit{ocultación de tokens} (\textit{token hiding}), inspirada en la arquitectura de credenciales anónimas LAASSO \cite{ma2026laasso}. Este mecanismo implementa un esquema de cifrado híbrido poscuántico que combina el algoritmo asimétrico sobre retículos modulares Kyber (ML-KEM) con el algoritmo de cifrado simétrico estándar AES \cite{ma2026laasso}.

Formalmente, el servidor central genera un par de claves de cifrado y descifrado $(vpk, vsk)$ mediante el algoritmo de generación de claves de Kyber \cite{ma2026laasso}:

\begin{equation}
(vpk, vsk) \leftarrow \text{Kyber.KeyGen}(1^\lambda)
\end{equation}

donde $vpk$ es la clave pública de verificación distribuida y $vsk$ es la clave privada custodiada en el \textit{Hardware Security Module} - HSM del servidor \cite{ma2026laasso}. Cuando el cliente construye la prueba ZKP no interactiva $\text{tok}'$, junto con el identificador único de ticket $\mathbf{t}$ y la marca de tiempo de validez $\nu$, ejecuta la fase de enmascaramiento $\text{SignOn}$ estructurada en los siguientes pasos \cite{ma2026laasso}:

\begin{enumerate}
    \item \textbf{Generación de Clave Simétrica Efímera:} El cliente muestrea de forma completamente aleatoria una clave simétrica de un solo uso $K \stackrel{\$}{\leftarrow} \{0, 1\}^\lambda$ (típicamente para AES-256) \cite{ma2026laasso}.
    \item \textbf{Encapsulamiento Asimétrico Poscuántico:} El cliente encripta la clave simétrica efímera $K$ utilizando la clave pública del servidor bancario $vpk$ mediante la función de cifrado de Kyber \cite{ma2026laasso}:
    \begin{equation}
    \text{tok}_1 = \text{Kyber.Enc}(vpk, K)
    \end{equation}
    \item \textbf{Cifrado Simétrico del Token de Autenticación:} El cliente encripta el paquete de datos que contiene la argumento de prueba $\text{tok}'$, el identificador $\mathbf{t}$ y el sello de tiempo $\nu$ empleando la clave $K$ a través de cifrado simétrico AES \cite{ma2026laasso}:
    \begin{equation}
    \text{tok}_2 = \text{AES.Enc}(K, \{\text{tok}', \mathbf{t}, \nu\})
    \end{equation}
    \item \textbf{Ensamblaje del Token Compuesto:} El token final transmitido hacia la red se define como la tupla $\text{tok} = (\text{tok}_1, \text{tok}_2)$ \cite{ma2026laasso}.
\end{enumerate}

\paragraph{Garantías de Seguridad y el Canal Ciego de Transmisión}

El procesamiento del token en la entidad receptora garantiza la verificación exclusiva por el destinatario autorizado. Al recibir el token compuesto $\text{tok} = (\text{tok}_1, \text{tok}_2)$, el servidor central recupera la clave simétrica $K = \text{Kyber.Dec}(vsk, \text{tok}_1)$ utilizando su clave privada $vsk$ y, posteriormente, descifra el contenido simétrico $\{\text{tok}', \mathbf{t}, \nu\} = \text{AES.Dec}(K, \text{tok}_2)$ para proceder con la verificación de la prueba ZKP sobre el retículo \cite{ma2026laasso}.

Las propiedades de seguridad del mecanismo de \textit{token hiding} se fundamentan en las siguientes garantías formales \cite{ma2026laasso}:

\begin{itemize}
    \item \textbf{Seguridad de Verificador Designado (\textit{Designated Verifier Security}):} Debido a la seguridad IND-CPA del esquema de cifrado Kyber sobre el problema Module-LWE, cualquier intermediario que no posea la clave privada $vsk$ es incapaz de descifrar $\text{tok}_1$ para extraer la clave $K$ \cite{ma2026laasso}.
    
    \item \textbf{Indistinguibilidad del Token frente a Ruido Uniforme:} Para un atacante que escuche el canal público, la combinación del encapsulamiento con Kyber y el cifrado con AES asegura que la estructura del token $\text{tok} = (\text{tok}_1, \text{tok}_2)$ sea computacionalmente indistinguible de secuencias de ruido aleatorio uniforme sobre el alfabeto binario $\{0, 1\}^*$ \cite{ma2026laasso}. La ventaja de un adversario $\mathcal{A}$ para distinguir entre tokens pertenecientes a dos transacciones distintas viene acotada por \cite{ma2026laasso}:
    \begin{equation}
    \left| \Pr[\mathcal{A}(\text{tok}_0) = 1] - \Pr[\mathcal{A}(\text{tok}_1) = 1] \right| \le \text{negl}(\lambda)
    \end{equation}
    
    \item \textbf{Paso a Canal Ciego de Transmisión (\textit{Blind Transmission Channel}):} El mediador queda relegado operacionalmente a actuar como un mero \textit{canal ciego de transmisión}. La terminal recibe el bloque cifrado $\text{tok}$, no puede leer su contenido ni evaluar si la prueba ZKP es válida o falsa, y se limita a reenviarlo al servidor central.
\end{itemize}

Esta ocultación profunda garantiza que el mediador no pueda almacenar credenciales ni perfilar al usuario, protegiendo de forma absoluta la privacidad de la comunicación y previniendo ataques de repetición mediante la verificación centralizada del identificador único de ticket $\mathbf{t}$ en el registro de consumo del servidor \cite{ma2026laasso}.


\section{Propuesta del Protocolo de Autenticación Poscuántica: PQLite-Auth v2}

\subsubsection{Arquitectura General y Modelo de Amenazas}

\paragraph{Definición Teórica y Objetivos del Protocolo}

Se propone \textit{PQLite-Auth v2}, un esquema de autenticación no interactiva de conocimiento cero (\textit{Non-Interactive Zero-Knowledge Argument of Knowledge} - NIZKAoK) optimizado para entornos de baja capacidad de cómputo y transmisión sobre canales no confiables \cite{yang2025multiserver, ma2026laasso}. El protocolo se fundamenta en la dureza del problema \textit{Module Learning with Errors} (MLWE) sobre el anillo polinomial $\mathcal{R}_q = \mathbb{Z}_q[x]/\langle x^{256} + 1 \rangle$, e integra un mecanismo de ocultación de tokens (\textit{Token Hiding}) y acuerdo de clave de sesión mediante reconciliación de errores en cuatro cuadrantes \cite{yang2025multiserver, ma2026laasso, lai2025leaky}.

El objetivo primario de PQLite-Auth v2 es permitir que un dispositivo cliente demuestre la posesión de una clave secreta de retículo $\mathbf{s} \in \mathcal{R}_q^k$ ante un servidor verificador designado, garantizando las siguientes propiedades criptográficas \cite{lombardi2022postquantum, ma2026laasso}:

\begin{itemize}
    \item \textbf{Resistencia Poscuántica y Solidez ($\text{Soundness}$):} Ningún adversario con acceso a computación clásica ni cuántica puede suplantar la identidad de un usuario legítimo sin resolver una instancia dura de MLWE \cite{regev2009lwe, badiezadegan2026complexity}.
    \item \textbf{Conocimiento Cero ($\text{Zero-Knowledge}$):} El servidor verificador y los nodos intermedios obtienen cero información estructural sobre la clave privada $\mathbf{s}$ durante las interacciones del protocolo \cite{goldwasser1989knowledge, lai2025leaky}.
    \item \textbf{Seguridad de Verificador Designado y No Vinculabilidad (\textit{Unlinkability}):} Las pruebas emitidas son computacionalmente indistinguibles de ruido aleatorio uniforme para cualquier entidad distinta del servidor central designado, impidiendo el rastreo o perfilamiento del cliente \cite{ma2026laasso, lombardi2022postquantum}.
    \item \textbf{Prevención de Ataques de Repetición (\textit{Anti-Replay Security}):} La inclusión de tickets de un solo uso procesados de forma determinista impide la reutilización de pruebas capturadas previamente en el canal público \cite{ma2026laasso}.
\end{itemize}

\paragraph{Entidades y Actores del Sistema}

La arquitectura de PQLite-Auth v2 se estructura formalmente a partir de tres actores abstractos con responsabilidades criptográficas delimitadas \cite{ma2026laasso, yang2025multiserver}:

\begin{enumerate}
    \item \textbf{Dispositivo Cliente / Elemento Seguro ($\mathcal{C}$):} Dispositivo de recursos limitados que almacena de forma inalterable en su memoria física protegida la clave secreta de retículo $\mathbf{s} \in \mathcal{R}_q^k$ y la clave pública de verificación del servidor $vpk$ \cite{mukil2026implementing, ma2026laasso}. El dispositivo $\mathcal{C}$ ejecuta la generación de la prueba NIZKAoK y el enmascaramiento híbrido sin exportar la clave secreta fuera del silicio \cite{mukil2026implementing}.
    \item \textbf{Terminal Intermedia / Nodo Relevador ($\mathcal{T}$):} Interfaz física o terminal de acceso potencialmente hostil que actúa como un \textit{canal ciego de transmisión} (\textit{blind relay}) \cite{ma2026laasso}. La terminal $\mathcal{T}$ reenvía los bloques cifrados emitidos por $\mathcal{C}$ hacia el servidor central sin capacidad de inspeccionar, descifrar o validar las pruebas criptográficas \cite{ma2026laasso}.
    \item \textbf{Servidor Verificador Designado ($\mathcal{S}$):} Entidad central de alta capacidad de cómputo o Módulo de Seguridad de Hardware (\textit{Hardware Security Module} - HSM) que custodia la clave privada de descifrado $vsk$ de Kyber, la base de datos de compromisos públicos $\mathbf{t} = \mathbf{A}\mathbf{s} + \mathbf{e} \pmod q$ y el registro de tickets consumidos (\textit{Spent Log}) \cite{yang2025multiserver, ma2026laasso}.
\end{enumerate}

\paragraph{Modelo de Amenazas}

El modelo de seguridad asume un adversario cuántico activo $\mathcal{A}$ que controla completamente el canal de comunicación entre $\mathcal{C}$, $\mathcal{T}$ y $\mathcal{S}$ \cite{lombardi2022postquantum, ma2026laasso}. Específicamente, se consideran los siguientes escenarios de ataque \cite{lai2025leaky, ma2026laasso}:

\begin{itemize}
    \item \textbf{Intercepción en Terminales Comprometidas:} El adversario $\mathcal{A}$ puede controlar la terminal intermedia $\mathcal{T}$ o instalar componentes maliciosos para capturar todo el tráfico transmitido por la red \cite{ma2026laasso}.
    \item \textbf{Consultas Semi-Adaptativas de Fuga (Leaky LWE):} El adversario $\mathcal{A}$ puede seleccionar matrices de desafío $\mathbf{L}$ o inyectar perturbaciones ruidosas para intentar medir componentes lineales de la clave secreta $\mathbf{s}$ a lo largo de múltiples transacciones \cite{lai2025leaky}.
    \item \textbf{Ataques Geométricos de Reducción de Bases (BKZ):} El adversario $\mathcal{A}$ dispone de algoritmos cuánticos de cribado (\textit{quantum sieving}) para intentar resolver el problema del vector más corto (SVP) sobre la base del retículo generado por $(\mathbf{A}, \mathbf{t})$ \cite{badiezadegan2026complexity, mukil2026implementing}.
\end{itemize}

\paragraph{Vector de Parámetros Concretos}

Para certificar el nivel de seguridad NIST Nivel 1 (128 bits de resistencia poscuántica) minimizando la sobrecarga de memoria en el dispositivo cliente, se fijan los parámetros en el anillo $\mathcal{R}_q = \mathbb{Z}_q[x]/\langle x^d + 1 \rangle$ \cite{yang2025multiserver, ma2026laasso}:

\begin{equation}
d = 256, \quad k = 3, \quad q = 3329, \quad \eta = 2
\end{equation}

La dimensión espacial total del retículo resultante es $N = k \times d = 3 \times 256 = 768$ \cite{yang2025multiserver}. El módulo primo de 16 bits $q = 3329$ admite la Transformada de Teoría de Números (NTT), permitiendo multiplicaciones polinomiales eficientes de complejidad $\mathcal{O}(d \log d)$ \cite{yang2025multiserver, mukil2026implementing}. La distribución de error corresponde a una Distribución Binomial Centrada $\text{CBD}(\eta=2)$ con desviación estándar $\sigma = 1$, implementada en tiempo constante $\mathcal{O}(1)$ \cite{yang2025multiserver}.

\paragraph{Secuencia Lógica Temporal del Protocolo}

El flujo operativo del protocolo PQLite-Auth v2 se ejecuta en cuatro fases secuenciales estrictas, coordinadas entre el dispositivo cliente $\mathcal{C}$ y el servidor verificador $\mathcal{S}$ a través del canal ciego $\mathcal{T}$ \cite{yang2025multiserver, ma2026laasso}:

\begin{enumerate}
    \item \textbf{Fase 1 (Generación de Pruebas y Ocultación en $\mathcal{C}$):} El cliente $\mathcal{C}$ muestrea ruido fresco, genera el argumento NIZKAoK determinista mediante Fiat-Shamir con Abortos, verifica la cota de norma, encapsula una clave simétrica efímera con Kyber ($vpk$) y encripta el token mediante AES-256 \cite{ma2026laasso}.
    \item \textbf{Fase 2 (Transmisión Ciega en $\mathcal{T}$):} La terminal $\mathcal{T}$ recibe el paquete cifrado $\text{tok} = (\text{tok}_1, \text{tok}_2)$ y lo reenvía al servidor $\mathcal{S}$ sin realizar procesamiento local \cite{ma2026laasso}.
    \item \textbf{Fase 3 (Verificación y Desencapsulamiento en $\mathcal{S}$):} El servidor $\mathcal{S}$ descifra el token usando su clave privada $vsk$, valida la frescura del ticket contra el \textit{Spent Log}, reconstruye el desafío y verifica las ecuaciones modulares y la cota de norma del retículo \cite{yang2025multiserver, ma2026laasso}.
    \item \textbf{Fase 4 (Acuerdo de Clave de Sesión $\mathcal{S} \to \mathcal{C}$):} Aprobada la autenticación, $\mathcal{S}$ deriva una pista de redondeo cruzado en 4 cuadrantes $\mathbf{v} \in \{0,1\}^{256}$ y la transmite a $\mathcal{C}$, permitiendo a ambas partes concordar la misma clave simétrica de 256 bits mediante reconciliación de Peikert \cite{yang2025multiserver}.
\end{enumerate}

\subsubsection{Generación del Argumento NIZKAoK y Ocultación de Tokens (Token Hiding)}

\paragraph{Fundamentación del Argumento NIZKAoK mediante Fiat-Shamir con Abortos}

La primera etapa operacional del protocolo PQLite-Auth v2 se ejecuta íntegramente dentro del entorno físicamente protegido del dispositivo cliente ($\mathcal{C}$) \cite{mukil2026implementing, ma2026laasso}. Su propósito es construir un argumento no interactivo de conocimiento de conocimiento cero (\textit{Non-Interactive Zero-Knowledge Argument of Knowledge} - NIZKAoK) que pruebe la posesión del secreto vectorial $\mathbf{s} \in \mathcal{R}_q^k$ asociado al compromiso público $\mathbf{t} = \mathbf{A}\mathbf{s} + \mathbf{e} \pmod q$, sin revelar componente alguna de $\mathbf{s}$ \cite{lyubashevsky2009fiat, ma2026laasso}.

Para evitar que la respuesta lineal filtre la geometría de la clave secreta, el protocolo adapta el paradigma de \textit{Fiat-Shamir con Abortos} (\textit{Fiat-Shamir with aborts}) ideado por Lyubashevsky \cite{lyubashevsky2009fiat, lyubashevsky2012lattice}. La generación de la prueba se ejecuta mediante la siguiente secuencia de operaciones algebraicas en el anillo $\mathcal{R}_q = \mathbb{Z}_q[x]/\langle x^{256} + 1 \rangle$ \cite{yang2025multiserver, ma2026laasso}:

\begin{enumerate}
    \item \textbf{Muestreo de Ruido Fresco de Enmascaramiento:} El dispositivo $\mathcal{C}$ genera de forma estocástica e independiente un vector de ruido fresco de enmascaramiento $\mathbf{y} \in \mathcal{R}_q^k$, cuyos coeficientes se muestrean a partir de la Distribución Binomial Centrada $\text{CBD}(\eta=2)$ con desviación estándar $\sigma = 1$ \cite{yang2025multiserver}:
    \begin{equation}
    \mathbf{y} \leftarrow \text{CBD}(\eta=2)^k \in \mathcal{R}_q^k
    \end{equation}
    
    \item \textbf{Cálculo del Compromiso Providencial ($\mathbf{w}$):} Utilizando la matriz pública global del sistema $\mathbf{A} \in \mathcal{R}_q^{k \times k}$, el cliente evalúa el vector de compromiso preliminar $\mathbf{w}$ \cite{yang2025multiserver, ma2026laasso}:
    \begin{equation}
    \mathbf{w} = \mathbf{A}\mathbf{y} \pmod q \in \mathcal{R}_q^k
    \end{equation}
    
    \item \textbf{Derivación Determinista del Desafío ($\mathbf{L}$):} Mediante la heurística de Fiat-Shamir en el modelo de oráculo aleatorio (\textit{Random Oracle Model} - ROM), $\mathcal{C}$ calcula la matriz de desafío no interactivo $\mathbf{L} \in \mathcal{R}_q^{k \times k}$ aplicando la función de hash criptográfica $H: \{0,1\}^* \to \mathcal{R}_q^{k \times k}$ sobre la concatenación del compromiso público $\mathbf{t}$, el compromiso fresco $\mathbf{w}$, un identificador único de ticket de un solo uso $\mathbf{t}_{\text{ticket}} \in \{0,1\}^{256}$ y la marca de tiempo de la sesión $\nu$ \cite{fiat1986how, lyubashevsky2009fiat, ma2026laasso}:
    \begin{equation}
    \mathbf{L} = H(\mathbf{A}, \mathbf{t}, \mathbf{w}, \mathbf{t}_{\text{ticket}}, \nu) \in \mathcal{R}_q^{k \times k}
    \end{equation}
    La inclusión explícita del par $(\mathbf{t}_{\text{ticket}}, \nu)$ en el oráculo garantiza la vinculación criptográfica entre la prueba NIZKAoK y la sesión temporal activa, previniendo la maleabilidad de los transcripts \cite{ma2026laasso}.
    
    \item \textbf{Evaluación del Vector de Respuesta ($\mathbf{z}$):} El cliente evalúa la combinación lineal modular entre la clave secreta $\mathbf{s}$, la matriz de desafío $\mathbf{L}$ y el ruido fresco $\mathbf{y}$ \cite{yang2025multiserver, ma2026laasso}:
    \begin{equation}
    \mathbf{z} = \mathbf{s}\mathbf{L} + \mathbf{y} \pmod q \in \mathcal{R}_q^k
    \end{equation}
    
    \item \textbf{Filtro de Muestreo de Rechazo (\textit{Rejection Sampling}):} Para romper la correlación estadística entre la respuesta $\mathbf{z}$ y el secreto $\mathbf{s}$, el algoritmo aplica una verificación de cota de norma infinita $\|\cdot\|_\infty$ \cite{lyubashevsky2009fiat, lyubashevsky2012lattice, lai2025leaky}:
    \begin{equation}
    \|\mathbf{z}\|_\infty \le \beta_z
    \end{equation}
    donde $\beta_z$ es el umbral de aceptación ajustado según el parámetro de escalamiento de Leaky LWE \cite{lai2025leaky}. Si $\|\mathbf{z}\|_\infty > \beta_z$, el procesador de $\mathcal{C}$ \textit{aborta la iteración de inmediato}, descarta la tupla $(\mathbf{y}, \mathbf{z}, \mathbf{w})$ y retorna de forma recursiva al Paso 1 para muestrear un nuevo vector de ruido $\mathbf{y}'$ \cite{lyubashevsky2009fiat, ma2026laasso}. Si $\|\mathbf{z}\|_\infty \le \beta_z$, la prueba NIZKAoK $(\mathbf{z}, \mathbf{w})$ se declara válida y el dispositivo avanza hacia la fase de enmascaramiento \cite{lyubashevsky2009fiat, ma2026laasso}.
\end{enumerate}

\paragraph{Mecanismo de Ocultación de Tokens (\textit{Token Hiding}) con Kyber y AES}

Aprobada la verificación de norma del argumento NIZKAoK en $\mathcal{C}$, el protocolo inicia el procedimiento de enmascaramiento para garantizar la propiedad de \textit{verificador designado} (\textit{designated verifier}) e impedir que la terminal intermedia $\mathcal{T}$ o cualquier observador del canal inspeccione el contenido de la prueba \cite{ma2026laasso}. El mecanismo de \textit{Token Hiding} implementa una envoltura de cifrado híbrido poscuántico estructurada en los siguientes pasos \cite{ma2026laasso}:

\begin{enumerate}
    \item \textbf{Generación de la Clave Simétrica Efímera ($K$):} El cliente $\mathcal{C}$ muestrea un bloque de aleatoriedad uniforme de 256 bits para instanciar una clave simétrica de un solo uso \cite{ma2026laasso}:
    \begin{equation}
    K \stackrel{\$}{\leftarrow} \{0, 1\}^{256}
    \end{equation}
    
    \item \textbf{Encapsulamiento Asimétrico Poscuántico ($\text{tok}_1$):} Utilizando la clave pública de descifrado del servidor central $vpk$ (correspondiente al estándar ML-KEM / Kyber-768 almacenado en la memoria inalterable de $\mathcal{C}$), el dispositivo encapsula la clave efímera $K$ \cite{ma2026laasso, mukil2026implementing}:
    \begin{equation}
    \text{tok}_1 = \text{Kyber.Enc}(vpk, K)
    \end{equation}
    El resultado $\text{tok}_1$ constituye un texto cifrado asimétrico sobre Module-LWE que solo puede ser revertido por la clave privada $vsk$ del servidor $\mathcal{S}$ \cite{ma2026laasso}.
    
    \item \textbf{Cifrado Simétrico AES-256 del Carga Útil ($\text{tok}_2$):} Utilizando la clave efímera $K$ bajo el algoritmo simétrico AES en modo GCM/CBC con vector de inicialización fresco, $\mathcal{C}$ encripta la tupla que contiene la prueba NIZKAoK, el compromiso $\mathbf{w}$, el identificador de ticket $\mathbf{t}_{\text{ticket}}$ y el sello de tiempo $\nu$ \cite{ma2026laasso}:
    \begin{equation}
    \text{tok}_2 = \text{AES.Enc}\left(K, \{\mathbf{z}, \mathbf{w}, \mathbf{t}_{\text{ticket}}, \nu\}\right)
    \end{equation}
    
    \item \textbf{Construcción y Emisión del Token Compuesto ($\text{tok}$):} El dispositivo cliente concatena las dos estructuras cifradas para conformar el token de autenticación final \cite{ma2026laasso}:
    \begin{equation}
    \text{tok} = (\text{tok}_1, \text{tok}_2)
    \end{equation}
    El paquete $\text{tok}$ se transmite desde $\mathcal{C}$ hacia la terminal intermedia $\mathcal{T}$ mediante la interfaz de comunicación \cite{ma2026laasso}.
\end{enumerate}

\paragraph{Transferencia Ciega en la Terminal Intermedia ($\mathcal{T}$)}

La terminal intermedia $\mathcal{T}$ recibe el paquete compuesto $\text{tok} = (\text{tok}_1, \text{tok}_2)$. Debido a la propiedad de indecodificabilidad bajo ataques de texto claro elegido (IND-CPA) de Kyber y AES-256, el bloque $\text{tok}$ resulta computacionalmente indistinguible de una secuencia de ruido aleatorio uniforme sobre el alfabeto binario para la terminal $\mathcal{T}$ \cite{ma2026laasso}:

\begin{equation}
\text{tok} \approx_c \mathcal{U}\left(\{0, 1\}^{|\text{tok}|}\right)
\end{equation}

Por consiguiente, la terminal $\mathcal{T}$ no puede descifrar la clave $K$, no puede extraer el vector de respuesta $\mathbf{z}$, ni puede verificar si las ecuaciones del retículo son válidas \cite{ma2026laasso}. Operacionalmente, $\mathcal{T}$ actúa como un \textit{canal ciego de retransmisión} (\textit{blind relay}), limitándose a empaquetar $\text{tok}$ dentro de la trama de red y reenviarlo hacia la dirección de red del servidor verificador designado $\mathcal{S}$ \cite{ma2026laasso}.

\subsubsection{Descifrado Centralizado, Verificación de Norma y Prevención de Replay}

\paragraph{Recepción y Descifrado Híbrido en el Servidor Designado}

El servidor central verificador ($\mathcal{S}$) es la única entidad que posee la clave privada de descifrado $vsk$ correspondiente al par de claves asimétricas de Kyber-768 instauradas en el sistema \cite{ma2026laasso, mukil2026implementing}. Al recibir el paquete cifrado $\text{tok} = (\text{tok}_1, \text{tok}_2)$ retransmitido por la terminal intermedia $\mathcal{T}$, el servidor $\mathcal{S}$ inicia el procedimiento de desencapsulamiento y descifrado simétrico estructurado en las siguientes etapas \cite{ma2026laasso}:

\begin{enumerate}
    \item \textbf{Recuperación de la Clave Simétrica Efímera ($K$):} Mediante el algoritmo de descifrado asimétrico de Kyber, el servidor $\mathcal{S}$ aplica su clave privada $vsk$ sobre la primera componente del token $\text{tok}_1$ \cite{ma2026laasso}:
    \begin{equation}
    K = \text{Kyber.Dec}(vsk, \text{tok}_1)
    \end{equation}
    La seguridad de la construcción garantiza que, salvo que el atacante resuelva una instancia dura del problema MLWE sobre la clave pública $vpk$, la clave simétrica de 256 bits $K$ se recupera de forma íntegra y sin alteraciones \cite{ma2026laasso}.

    \item \textbf{Descifrado Simétrico de la Carga Útil ($\text{payload}$):} Utilizando la clave efímera $K$ recuperada, $\mathcal{S}$ ejecuta el descifrado AES-256 sobre la segunda componente $\text{tok}_2$ para extraer la carga útil en texto claro \cite{ma2026laasso}:
    \begin{equation}
    \{\mathbf{z}, \mathbf{w}, \mathbf{t}_{\text{ticket}}, \nu\} = \text{AES.Dec}(K, \text{tok}_2)
    \end{equation}
    donde $\mathbf{z} \in \mathcal{R}_q^k$ es la respuesta NIZKAoK, $\mathbf{w} \in \mathcal{R}_q^k$ es el vector de compromiso fresco, $\mathbf{t}_{\text{ticket}} \in \{0,1\}^{256}$ es el identificador único de ticket y $\nu$ es el sello temporal de la sesión \cite{yang2025multiserver, ma2026laasso}.
\end{enumerate}

\paragraph{Verificación Anti-Replay y Frescura Temporal mediante Registro de Tickets}

Antes de procesar la prueba sobre el retículo, el servidor $\mathcal{S}$ valida las propiedades de no maleabilidad y frescura temporal del paquete para neutralizar ataques de repetición (\textit{replay attacks}) e intercepción diferida \cite{lombardi2022postquantum, ma2026laasso}. Esta validación se rige por las siguientes condiciones de control \cite{ma2026laasso}:

\begin{enumerate}
    \item \textbf{Control de Ventana Temporal ($\Delta t$):} El servidor compara la marca de tiempo $\nu$ contenida en la carga útil descifrada contra el reloj interno del servidor $\nu_{\text{actual}}$. Si $|\nu_{\text{actual}} - \nu| > \Delta t$ (donde $\Delta t$ define la ventana de tolerancia máxima del sistema), el servidor aborta la transacción por caducidad temporal \cite{ma2026laasso}.
    
    \item \textbf{Verificación en el Registro de Consumo (\textit{Spent Log}):} El servidor $\mathcal{S}$ realiza una consulta de búsqueda determinista sobre el registro centralizado de tickets consumidos $\mathcal{L}_{\text{spent}}$ empleando el identificador $\mathbf{t}_{\text{ticket}}$ \cite{ma2026laasso}:
    \begin{equation}
    \mathbf{t}_{\text{ticket}} \stackrel{?}{\in} \mathcal{L}_{\text{spent}}
    \end{equation}
    Si $\mathbf{t}_{\text{ticket}} \in \mathcal{L}_{\text{spent}}$, se detecta un intento de reutilización de un token capturado previamente en el canal público, por lo que el servidor $\mathcal{S}$ rechaza de inmediato la solicitud y emite una alerta de seguridad \cite{ma2026laasso}. Si $\mathbf{t}_{\text{ticket}} \notin \mathcal{L}_{\text{spent}}$, el servidor añade formalmente el identificador al registro ($\mathcal{L}_{\text{spent}} \leftarrow \mathcal{L}_{\text{spent}} \cup \{\mathbf{t}_{\text{ticket}}\}$) para invalidar cualquier retransmisión futura de dicho ticket \cite{ma2026laasso}.
\end{enumerate}

\paragraph{Reconstrucción del Desafío y Verificación Geométrica sobre el Retículo}

Superado el filtro anti-replay, el servidor $\mathcal{S}$ valida la solidez matemática de la prueba NIZKAoK evaluando la relación algebraica entre la respuesta $\mathbf{z}$, el compromiso público $\mathbf{t} = \mathbf{A}\mathbf{s} + \mathbf{e} \pmod q$ almacenado en la base de datos central y el compromiso providencial $\mathbf{w}$ \cite{yang2025multiserver, ma2026laasso}:

\begin{enumerate}
    \item \textbf{Reconstrucción del Desafío Fiat-Shamir ($\mathbf{L}$):} El servidor $\mathcal{S}$ invoca la función de hash oráculo $H$ sobre los elementos recuperados para reconstruir exactamente la matriz de desafío no interactiva $\mathbf{L} \in \mathcal{R}_q^{k \times k}$ \cite{fiat1986how, lyubashevsky2009fiat, ma2026laasso}:
    \begin{equation}
    \mathbf{L} = H(\mathbf{A}, \mathbf{t}, \mathbf{w}, \mathbf{t}_{\text{ticket}}, \nu) \in \mathcal{R}_q^{k \times k}
    \end{equation}

    \item \textbf{Comprobación Estricta de Cota de Norma ($\|\mathbf{z}\|_\infty$):} El servidor verifica que el vector de respuesta $\mathbf{z}$ cumpla la restricción de norma máxima \cite{lyubashevsky2009fiat, lai2025leaky}:
    \begin{equation}
    \|\mathbf{z}\|_\infty \stackrel{?}{\le} \beta_z
    \end{equation}
    Si $\|\mathbf{z}\|_\infty > \beta_z$, el servidor rechaza la prueba de inmediato. Esta verificación garantiza que un adversario no pueda presentar vectores de respuesta arbitrariamente largos para satisfacer la ecuación modular sin conocer el secreto corto $\mathbf{s}$ \cite{lyubashevsky2009fiat, lai2025leaky}.

    \item \textbf{Evaluación de Consistencia Algebraica sobre MLWE:} El servidor $\mathcal{S}$ recupera la matriz pública $\mathbf{A}$ y el compromiso $\mathbf{t}$ asociados a la identidad del cliente, y evalúa el producto matricial modular \cite{yang2025multiserver, ma2026laasso}:
    \begin{equation}
    \mathbf{A}\mathbf{z} - \mathbf{t}\mathbf{L} \pmod q \in \mathcal{R}_q^k
    \end{equation}
    Sustituyendo las definiciones del probador ($\mathbf{z} = \mathbf{s}\mathbf{L} + \mathbf{y}$ y $\mathbf{t} = \mathbf{A}\mathbf{s} + \mathbf{e}$), se observa la cancelación del término de secreto principal \cite{yang2025multiserver}:
    \begin{align}
    \mathbf{A}\mathbf{z} - \mathbf{t}\mathbf{L} &\equiv \mathbf{A}(\mathbf{s}\mathbf{L} + \mathbf{y}) - (\mathbf{A}\mathbf{s} + \mathbf{e})\mathbf{L} \pmod q \nonumber \\
    &\equiv \mathbf{A}\mathbf{s}\mathbf{L} + \mathbf{A}\mathbf{y} - \mathbf{A}\mathbf{s}\mathbf{L} - \mathbf{e}\mathbf{L} \pmod q \nonumber \\
    &\equiv \mathbf{A}\mathbf{y} - \mathbf{e}\mathbf{L} \pmod q \nonumber \\
    &\equiv \mathbf{w} - \mathbf{e}\mathbf{L} \pmod q
    \end{align}

    \item \textbf{Validación de Cota de Ruido Residual:} Dado que $\mathbf{w} = \mathbf{A}\mathbf{y} \pmod q$, la diferencia entre la expresión calculada $\mathbf{A}\mathbf{z} - \mathbf{t}\mathbf{L}$ y el compromiso providencial $\mathbf{w}$ viene dada por el término de error perturbador $-\mathbf{e}\mathbf{L}$ \cite{yang2025multiserver}. El servidor verifica que la norma del error residual esté acotada por la cota teórica $\beta_e$ \cite{yang2025multiserver, lai2025leaky}:
    \begin{equation}
    \|\mathbf{A}\mathbf{z} - \mathbf{t}\mathbf{L} - \mathbf{w}\|_\infty \stackrel{?}{\le} \beta_e
    \end{equation}
    Si la diferencia respeta la cota $\beta_e$, el servidor declara la prueba NIZKAoK como \textit{válida} y confirma la autenticación del dispositivo cliente $\mathcal{C}$, habilitando la ejecución de la fase final de acuerdo de clave de sesión \cite{yang2025multiserver, ma2026laasso}.
\end{enumerate}

\subsubsection{Acuerdo de Clave de Sesión mediante Reconciliación en Cuatro Cuadrantes}

\paragraph{Fundamentación de la Reconciliación de Errores sobre Retículos}

Tras la verificación exitosa de la prueba NIZKAoK y la confirmación de autenticidad en el servidor ($\mathcal{S}$), el protocolo culmina con la generación de una clave simétrica de sesión compartida $K_{\text{sesión}} \in \{0, 1\}^{256}$ entre el dispositivo cliente ($\mathcal{C}$) y el servidor ($\mathcal{S}$) \cite{yang2025multiserver, rewal2023quantumsafe}. Para evitar el uso de mecanismos adicionales de cifrado asimétrico que incrementarían la sobrecarga computacional del dispositivo embebido, ambas partes derivan la clave mediante la técnica de \textit{reconciliación de errores de Peikert} sobre el anillo $\mathcal{R}_q = \mathbb{Z}_q[x]/\langle x^{256} + 1 \rangle$ con módulo $q = 3329$ \cite{yang2025multiserver, rewal2023quantumsafe}.

En un protocolo sobre retículos, las multiplicaciones algebraicas independientes ejecutadas por el cliente y el servidor producen dos vectores continuos ruidosos $\mathbf{w}_A, \mathbf{w}_B \in \mathcal{R}_q$ que difieren únicamente por una pequeña perturbación de error LWE acumulado \cite{yang2025multiserver, rewal2023quantumsafe}:

\begin{equation}
\mathbf{w}_A \approx \mathbf{w}_B \pmod q \implies \mathbf{w}_A = \mathbf{w}_B + \mathbf{e}_{\text{rec}} \pmod q
\end{equation}

Debido a la naturaleza discreta de los bordes de redondeo modular en $\mathbb{Z}_q$, la aplicación de una función de redondeo directa sobre coeficientes cercanos a las fronteras provocaría discrepancias de bits si el error $\mathbf{e}_{\text{rec}}$ hace saltar el valor al cuadrante contiguo \cite{yang2025multiserver}. La reconciliación en cuatro cuadrantes resuelve el problema de la frontera mediante la emisión de un vector de señal auxiliar $\mathbf{v} \in \{0,1\}^{256}$ transmitido en texto claro, el cual orienta al cliente para corregir la zona de decisión sin filtrar información sobre los bits de la clave final \cite{yang2025multiserver}.

\paragraph{Partición Modular y Funciones de Peikert en $\mathbb{Z}_{3329}$}

El espacio del anillo modular $\mathbb{Z}_q$ con $q = 3329$ se particiona formalmente en cuatro cuadrantes de tamaño aproximado $\lfloor q/4 \rfloor \approx 832$ unidades \cite{yang2025multiserver}:

\begin{itemize}
    \item \textbf{Cuadrante 0 ($Q_0$):} $w \in \left[0, \lfloor q/4 \rfloor\right) \implies \text{Bit de Clave } K = 0, \quad \text{Señal } v = 0 \text{ (Par)}$
    \item \textbf{Cuadrante 1 ($Q_1$):} $w \in \left[\lfloor q/4 \rfloor, \lfloor q/2 \rfloor\right) \implies \text{Bit de Clave } K = 0, \quad \text{Señal } v = 1 \text{ (Impar)}$
    \item \textbf{Cuadrante 2 ($Q_2$):} $w \in \left[\lfloor q/2 \rfloor, \lfloor 3q/4 \rfloor\right) \implies \text{Bit de Clave } K = 1, \quad \text{Señal } v = 0 \text{ (Par)}$
    \item \textbf{Cuadrante 3 ($Q_3$):} $w \in \left[\lfloor 3q/4 \rfloor, q\right) \implies \text{Bit de Clave } K = 1, \quad \text{Señal } v = 1 \text{ (Impar)}$
\end{itemize}

Las tres funciones algebraicas fundamentales de Peikert adaptadas al módulo $q = 3329$ se definen coeficiente a coeficiente sobre $w \in \mathbb{Z}_q$ de la siguiente manera \cite{yang2025multiserver}:

\begin{enumerate}
    \item \textbf{Función de Redondeo Modular (Extracción de Bit de Clave):} Extrae el bit principal según la mitad del círculo modular \cite{yang2025multiserver}:
    \begin{equation}
    \lfloor w \rfloor_{q,2} = \left\lfloor \frac{2w}{q} \right\rceil \bmod 2 \in \{0, 1\}
    \end{equation}
    
    \item \textbf{Función de Redondeo Cruzado (Generación de Señal Auxiliary $v$):} Identifica si el valor pertenece a un cuadrante Par ($Q_0, Q_2$) o Impar ($Q_1, Q_3$) \cite{yang2025multiserver}:
    \begin{equation}
    \langle w \rangle_{q,2} = \left\lfloor \frac{4w}{q} \right\rfloor \bmod 2 \in \{0, 1\}
    \end{equation}
    
    \item \textbf{Función de Reconciliación de Errores ($\text{rec}$):} Reconstruye el bit de clave idéntico al del servidor utilizando el valor local del cliente $w_A$ y la pista de señal $v$ \cite{yang2025multiserver}:
    \begin{equation}
    \text{rec}(w_A, v) = \left\lfloor \frac{2}{q} \cdot \left( w_A - \left( (2v + 1) \cdot \frac{q}{4} \right) \right) \right\rceil \bmod 2 \in \{0, 1\}
    \end{equation}
\end{enumerate}

\paragraph{Secuencia de Generación de la Clave de Sesión}

El procedimiento de concordancia de clave simétrica entre $\mathcal{S}$ y $\mathcal{C}$ se ejecuta mediante los siguientes pasos numerados \cite{yang2025multiserver, rewal2023quantumsafe}:

\begin{enumerate}
    \item \textbf{Cálculo del Polinomio Base en el Servidor ($\mathbf{w}_B$):} El servidor $\mathcal{S}$ combina el vector de respuesta del cliente con su secreto local para evaluar el polinomio de retículo $\mathbf{w}_B \in \mathcal{R}_q$ de $d = 256$ coeficientes \cite{yang2025multiserver}:
    \begin{equation}
    \mathbf{w}_B = \mathbf{b}_{\mathcal{C}}^T \mathbf{s}_{\mathcal{S}} + \mathbf{e}_{\mathcal{S}} \pmod q \in \mathcal{R}_q
    \end{equation}
    
    \item \textbf{Extracción Central de la Clave y Derivación de la Pista ($\mathbf{v}$):} El servidor $\mathcal{S}$ deriva simultáneamente la clave simétrica del servidor $K_{\mathcal{S}} \in \{0, 1\}^{256}$ y el vector de señal auxiliar $\mathbf{v} \in \{0, 1\}^{256}$ aplicando las funciones de Peikert coeficiente a coeficiente sobre los 256 elementos de $\mathbf{w}_B$ \cite{yang2025multiserver}:
    \begin{equation}
    K_{\mathcal{S}}^{(i)} = \lfloor \mathbf{w}_B^{(i)} \rfloor_{q,2}, \quad \mathbf{v}^{(i)} = \langle \mathbf{w}_B^{(i)} \rangle_{q,2} \quad \text{para } i \in \{0, \dots, 255\}
    \end{equation}
    El vector de señal $\mathbf{v}$ posee una longitud de 256 bits (exactamente 32 bytes de carga útil) \cite{yang2025multiserver}.
    
    \item \textbf{Transmisión de la Pista en Texto Claro ($\mathbf{v}$):} El servidor $\mathcal{S}$ emite el bloque de 32 bytes $\mathbf{v}$ a través del canal de red. La terminal intermedia $\mathcal{T}$ retransmite $\mathbf{v}$ de forma transparente hacia el dispositivo cliente $\mathcal{C}$ \cite{yang2025multiserver}.
    
    \item \textbf{Reconciliación Paralela en el Cliente ($\mathcal{C}$):} El cliente $\mathcal{C}$ recupera $\mathbf{v}$, calcula su valor local $\mathbf{w}_A = \mathbf{b}_{\mathcal{S}}^T \mathbf{s}_{\mathcal{C}} \pmod q \in \mathcal{R}_q$ y ejecuta la función de reconciliación en paralelo sobre sus 256 coeficientes \cite{yang2025multiserver}:
    \begin{equation}
    K_{\mathcal{C}}^{(i)} = \text{rec}\left(\mathbf{w}_A^{(i)}, \mathbf{v}^{(i)}\right) \quad \text{para } i \in \{0, \dots, 255\}
    \end{equation}
    Si la norma del error residual respeta la cota $\|\mathbf{e}_{\text{rec}}\|_\infty < q/8 \approx 416$, la propiedad de convergencia del teorema de Peikert garantiza la coincidencia idéntica de la clave simétrica de 256 bits \cite{yang2025multiserver, rewal2023quantumsafe}:
    \begin{equation}
    K_{\mathcal{C}} = K_{\mathcal{S}} = K_{\text{sesión}} \in \{0, 1\}^{256}
    \end{equation}
\end{enumerate}

\paragraph{Inmunidad de la Pista $\mathbf{v}$ en el Canal Público}

La transmisión del vector de señal $\mathbf{v}$ en texto claro a través del canal no compromete la seguridad de la clave de sesión $K_{\text{sesión}}$ frente a un adversario $\mathcal{A}$ que intercepte el tráfico \cite{yang2025multiserver}. Debido a la simetría geométrica de la partición en cuatro cuadrantes, una señal $v^{(i)} = 0$ indica que el coeficiente está en un cuadrante Par, el cual corresponde con idéntica probabilidad del 50\% tanto al Cuadrante 0 ($Q_0 \implies \text{Bit } 0$) como al Cuadrante 2 ($Q_2 \implies \text{Bit } 1$) \cite{yang2025multiserver}:

\begin{equation}
\Pr\left(K^{(i)} = 0 \mid v^{(i)} = 0\right) = \Pr\left(K^{(i)} = 1 \mid v^{(i)} = 0\right) = \frac{1}{2}
\end{equation}

Por consiguiente, la distribución de la pista $\mathbf{v}$ es estadísticamente independiente del vector de clave $K_{\text{sesión}}$, garantizando que el adversario obtenga cero bits de ventaja computacional sobre la clave simétrica acordada \cite{yang2025multiserver, rewal2023quantumsafe}.

\subsubsection{Análisis Teórico de Seguridad, Resistencia Poscuántica y Complejidad}

\paragraph{Sustento Criptográfico mediante Leaky LWE y Módulos Compactos}

La seguridad formal del protocolo PQLite-Auth v2 ante la reducción del módulo a $q = 3329$ se fundamenta en la teoría de \textit{Leaky LWE} formalizada por Lai, Swarnakar y Woo \cite{lai2025leaky}. En las construcciones ZKP tradicionales basadas en retículos, la demostración de la propiedad de conocimiento cero requiere la técnica de \textit{inundación de ruido} (\textit{noise flooding}), la cual fuerza a escalar el módulo a enteros gigantes ($q > 2^{128}$) y las dimensiones a valores masivos para evitar que las respuestas filtren la clave secreta \cite{lai2025leaky}.

Bajo el marco de Leaky LWE, se demuestra que la presencia de fugas lineales ruidosas de la forma $\mathbf{l}^T = (\mathbf{s}^T, \mathbf{e}^T)\mathbf{L} + \mathbf{f}^T \pmod q$ no reduce la dureza computacional del problema LWE subyacente, siempre que la matriz de desafío $\mathbf{L}$ satisfaga la cota de norma espectral $\|\mathbf{L}^T \mathbf{L}\| \le \beta_L$ \cite{lai2025leaky}. Dado que la desviación del ruido de fuga $\sigma_f$ únicamente escala en función de la raíz cuadrada del número de consultas $Q$ ($q \ge c \cdot \sqrt{Q}$), el protocolo opera de forma probadamente segura con el módulo compacto $q = 3329$ (16 bits) \cite{yang2025multiserver, lai2025leaky}. Las respuestas emitidas resultan computacionalmente indistinguibles de ruido aleatorio uniforme sobre $\mathbb{Z}_q$, garantizando cero fuga de información sobre la clave privada $\mathbf{s}$ \cite{lai2025leaky}.

\paragraph{Análisis de Dureza Geométrica y Resistencia ante Ataques BKZ}

La evaluación de la resistencia poscuántica concreta del protocolo se determina mediante la estimación de la complejidad del algoritmo de reducción de bases de retículos BKZ (\textit{Block-Korkine-Zolotarev}) ejecutado sobre la instancia Module-LWE del sistema \cite{badiezadegan2026complexity, mukil2026implementing}.

En el esquema PQLite-Auth v2, la matriz de polinomios posee rango $k = 3$ sobre el anillo $\mathcal{R}_q$ de grado $d = 256$ \cite{yang2025multiserver}. Aunque la dimensión matricial sea $k=3$, la dimensión espacial total del retículo euclídeo proyectado es \cite{yang2025multiserver}:
\begin{equation}
N = k \times d = 3 \times 256 = 768
\end{equation}

Para recuperar la clave corta $(\mathbf{s}, \mathbf{e})$ en un retículo de dimensión $N = 768$ con módulo $q = 3329$ y desviación de ruido $\sigma = 1$, un atacante cuantitativo debe reducir la base del retículo dual hasta alcanzar un factor de Hermite raíz $\delta_0 \approx 1.0039$ \cite{badiezadegan2026complexity}. La heurística de Gauss para el perfil de BKZ establece que lograr un factor de calidad $\delta_0 \le 1.0039$ exige fijar un tamaño de bloque de reducción de subespacio $\beta_{\text{BKZ}} \ge 400$ \cite{badiezadegan2026complexity}.

La complejidad temporal de resolver el problema del vector más corto (SVP) exacto dentro de cada sub-bloque de dimensión $\beta_{\text{BKZ}}$ utilizando los algoritmos más eficientes de cribado cuántico (\textit{quantum sieving}) se modela mediante la cota \cite{badiezadegan2026complexity, mukil2026implementing}:
\begin{equation}
T_{\text{Sieve}}(\beta_{\text{BKZ}}) \approx 2^{0.265 \cdot \beta_{\text{BKZ}}} \text{ operaciones cuánticas}
\end{equation}

Sustituyendo el tamaño de bloque mínimo $\beta_{\text{BKZ}} = 400$ e integrando la sobrecarga del número de recorridos (\textit{tours}) e iteraciones de la ventana sobre las 768 dimensiones, la complejidad computacional total del ataque BKZ resulta en \cite{badiezadegan2026complexity, mukil2026implementing}:
\begin{equation}
T_{\text{BKZ}}(400) \ge 2^{160} \text{ operaciones cuánticas}
\end{equation}

Puesto que $T_{\text{BKZ}} \ge 2^{160}$ supera con holgura la barrera de $2^{128}$ operaciones exigida por el estándar NIST Nivel 1 (equivalente a la seguridad de CRYSTALS-Kyber 768), el protocolo se certifica como formalmente inexpugnable ante ataques geométricos y reducciones duales de bases \cite{badiezadegan2026complexity, mukil2026implementing}.

\paragraph{Sobrecarga de Memoria, Comunicación y Cómputo en el Dispositivo Cliente}

El diseño de PQLite-Auth v2 optimiza la eficiencia de recursos en el dispositivo cliente ($\mathcal{C}$) en tres ejes críticos \cite{mukil2026implementing, ma2026laasso}:

\begin{enumerate}
    \item \textbf{Huella de Memoria Estática y Dinámica (\textit{Memory Footprint}):}
    \begin{itemize}
        \item \textit{Clave Privada ($\mathbf{s}$):} Vector de $k=3$ polinomios de 256 coeficientes en $\mathbb{Z}_{3329}$ (almacenado de forma compacta utilizando 12 bits por coeficiente): $3 \times 256 \times 12 / 8 = 1.152 \text{ bytes} \approx 1.15 \text{ KB}$ \cite{yang2025multiserver}.
        \item \textit{Matriz Pública ($\mathbf{A}$):} No se almacena explícitamente en la memoria Flash del cliente; se expande de forma determinista desde una semilla uniforme de 32 bytes ($\rho$) mediante la función XOF SHAKE-128, reduciendo el consumo de almacenamiento estático \cite{yang2025multiserver}.
        \item \textit{Clave de Verificación del Servidor ($vpk$):} Clave pública Kyber-768 almacenada en memoria de solo lectura: $1.184 \text{ bytes} \approx 1.18 \text{ KB}$ \cite{mukil2026implementing}.
    \end{itemize}

    \item \textbf{Ancho de Banda y Sobrecarga de Comunicación (\textit{Communication Overhead}):}
    \begin{itemize}
        \item \textit{Token de Autenticación ($\text{tok} = (\text{tok}_1, \text{tok}_2)$):} El texto cifrado asimétrico $\text{tok}_1$ (Kyber-768) requiere $1.088 \text{ bytes}$. La carga útil simétrica $\text{tok}_2$ (AES-256) que contiene $(\mathbf{z}, \mathbf{w}, \mathbf{t}_{\text{ticket}}, \nu)$ requiere aproximadamente $2.304 \text{ bytes}$. La transmisión total emitida por $\mathcal{C}$ es de $\approx 3.39 \text{ KB}$ \cite{ma2026laasso}.
        \item \textit{Pista de Reconciliación ($\mathbf{v}$):} La respuesta enviada por el servidor para el acuerdo de clave de sesión consiste en un vector binario de 256 bits, lo que equivale a solo $32 \text{ bytes}$ de tráfico de entrada en el canal \cite{yang2025multiserver}.
    \end{itemize}

    \item \textbf{Eficiencia Computacional y Mitigación de Canales Laterales:}
    \begin{itemize}
        \item \textit{Multiplicación Polinomial NTT:} La transformada NTT sobre $\mathbb{Z}_{3329}[x]/\langle x^{256} + 1 \rangle$ reduce la complejidad de multiplicación de $\mathcal{O}(d^2)$ a $\mathcal{O}(d \log d)$, requiriendo menos de 10.000 ciclos de reloj por producto en microcontroladores de 32 bits \cite{yang2025multiserver, mukil2026implementing}.
        \item \textit{CBD en Tiempo Constante:} El muestreo de la Distribución Binomial Centrada $\text{CBD}(\eta=2)$ se ejecuta mediante operaciones lógicas a nivel de palabra en tiempo constante $\mathcal{O}(1)$, garantizando inmunidad frente a ataques de análisis de tiempo y potencia por canales laterales (\textit{Side-Channel Attacks} - SCA) \cite{yang2025multiserver, mukil2026implementing}.
    \end{itemize}
\end{enumerate}


\section{Evaluación experimental y benchmarking de protocolos de autenticación}

\label{chap:benchmark}

Este capítulo presenta la metodología experimental utilizada para evaluar y comparar
las diferentes alternativas criptográficas consideradas en este trabajo. El objetivo
no consiste únicamente en determinar cuál implementación presenta el menor tiempo de
ejecución, sino en estudiar los compromisos existentes entre rendimiento, consumo de
recursos, comunicación, seguridad, complejidad de implementación y requisitos
operacionales.

% ============================================================
\subsection{Objetivos de la evaluación}
\label{sec:benchmark-objectives}

La evaluación experimental persigue los siguientes objetivos:

\begin{enumerate}
    \item Comparar el costo computacional de los diferentes protocolos durante las
    principales etapas de autenticación.

    \item Comparar el consumo de memoria y almacenamiento requerido por cada
    alternativa.

    \item Determinar el volumen de información que debe transmitirse durante una
    autenticación y estudiar su comportamiento ante diferentes condiciones de red.

    \item Analizar las condiciones de seguridad asociadas con los parámetros utilizados
    por cada protocolo.

    \item Examinar los principales ataques relevantes para cada familia criptográfica
    y las condiciones necesarias para mitigarlos.

    \item Comparar la complejidad de implementación, configuración y despliegue de cada
    alternativa.

    \item Analizar las dependencias de confianza y los requisitos de infraestructura de
    cada protocolo.

    \item Evaluar el impacto que tendría cada alternativa sobre la experiencia de un
    usuario final.

    \item Identificar los principales compromisos entre seguridad, rendimiento,
    comunicación, memoria y facilidad de despliegue.

    \item Comparar el protocolo propuesto en este trabajo con las demás alternativas
    bajo las mismas condiciones experimentales.
\end{enumerate}


% ============================================================
\subsection{Principio de comparabilidad}
\label{sec:benchmark-comparability}

Para garantizar una comparación significativa, todos los candidatos implementan una
misma funcionalidad abstracta de autenticación.

El flujo experimental se divide conceptualmente en dos momentos: registro
(\textit{enrollment}) y autenticación.

\subsubsection{Registro}

Durante el registro se ejecutan las operaciones necesarias para preparar al usuario y
al verificador antes de realizar una autenticación.

De manera abstracta, esta etapa se representa como:

\begin{equation}
    (pk, sk, \theta) \leftarrow \operatorname{Setup}()
\end{equation}

donde:

\begin{itemize}
    \item $pk$ representa la información pública necesaria para verificar al usuario;
    \item $sk$ representa el material secreto conservado por el usuario o demostrador;
    \item $\theta$ representa parámetros adicionales requeridos por el protocolo.
\end{itemize}

Dependiendo de la construcción, esta operación puede incluir generación de claves,
muestreo de vectores, construcción de matrices, generación de parámetros públicos o
precomputaciones.

\subsubsection{Autenticación}

Una autenticación se modela mediante las siguientes operaciones:

\begin{align}
    c &\leftarrow \operatorname{Challenge}(pk), \\
    r &\leftarrow \operatorname{Respond}(sk,c), \\
    b &\leftarrow \operatorname{Verify}(pk,c,r),
\end{align}

donde:

\begin{itemize}
    \item $c$ es el reto generado por el verificador;
    \item $r$ es la respuesta generada por el usuario;
    \item $b \in \{0,1\}$ representa el resultado de la verificación.
\end{itemize}

Para una ejecución válida se espera que:

\begin{equation}
    \operatorname{Verify}(pk,c,
    \operatorname{Respond}(sk,c)) = 1.
\end{equation}

Esta abstracción permite medir las diferentes construcciones utilizando las mismas
etapas funcionales aun cuando las operaciones matemáticas internas sean diferentes.

\subsubsection{Separación entre registro y autenticación}

El tiempo de generación de claves no se incorpora automáticamente al tiempo de cada
autenticación, puesto que normalmente corresponde a una operación realizada durante
el registro inicial.

Por lo tanto, se reportan por separado:

\begin{equation}
    T_{\text{enrollment}}
\end{equation}

y

\begin{equation}
    T_{\text{authentication}}
    =
    T_{\text{challenge}}
    +
    T_{\text{response}}
    +
    T_{\text{verification}}.
\end{equation}

Cuando una implementación requiera operaciones adicionales en alguna de estas etapas,
estas deberán documentarse y contabilizarse explícitamente.


% ============================================================
\subsection{Protocolos evaluados}
\label{sec:benchmark-protocols}

La Tabla~\ref{tab:protocol-summary} resume las alternativas estudiadas.

\begin{table}[H]
\centering
\caption{Protocolos incluidos en la evaluación experimental.}
\label{tab:protocol-summary}
\small
\begin{tabularx}{\textwidth}{lYYY}
\toprule
\textbf{Protocolo} &
\textbf{Fundamento matemático} &
\textbf{Tipo de construcción} &
\textbf{Rol en la comparación} \\
\midrule

ECDSA &
Problema del logaritmo discreto sobre curvas elípticas &
Criptografía tradicional &
Referencia de industria \\

LWE estándar &
Learning With Errors &
Retículos &
Candidato post-cuántico \\

Binary-LWE &
Learning With Errors con secreto restringido &
Retículos &
Variante LWE \\

Ring-LWE &
Learning With Errors sobre estructuras algebraicas de anillos &
Retículos estructurados &
Variante estructurada \\

LWR &
Learning With Rounding &
Retículos &
Variante basada en redondeo \\

LWE propuesto &
\todoresult{describir fundamento exacto} &
Retículos &
Propuesta de este trabajo \\

\bottomrule
\end{tabularx}
\end{table}

ECDSA se utiliza como referencia de una alternativa ampliamente adoptada en sistemas
criptográficos tradicionales. Las construcciones basadas en LWE y problemas
relacionados permiten estudiar alternativas fundamentadas en retículos y analizar sus
propiedades frente a los requisitos considerados en este proyecto.

La comparación se realiza utilizando un mismo modelo funcional de autenticación y un
entorno experimental controlado. De esta manera, las diferencias reportadas buscan
reflejar las características de cada alternativa criptográfica y no diferencias
artificiales introducidas por procedimientos de medición distintos.

\subsection{Entorno experimental}
\label{sec:experimental-environment}

Todos los experimentos fueron ejecutados bajo el mismo en   torno de hardware y software.
La Tabla~\ref{tab:environment} documenta la configuración utilizada con el objetivo de
facilitar la reproducibilidad de los resultados.

\begin{table}[H]
\centering
\caption{Entorno utilizado para ejecutar los experimentos.}
\label{tab:environment}
\begin{tabularx}{0.9\textwidth}{lY}
\toprule
\textbf{Elemento} & \textbf{Configuración} \\
\midrule

Procesador &
\todoresult{modelo exacto de CPU} \\

Arquitectura &
\todoresult{x86\_64, ARM64, etc.} \\

Núcleos físicos &
\todoresult{valor} \\

Núcleos lógicos &
\todoresult{valor} \\

Memoria RAM &
\todoresult{capacidad total} \\

Sistema operativo &
\todoresult{distribución y versión} \\

Kernel &
\todoresult{versión} \\

Docker &
\todoresult{versión} \\

Imagen base &
\todoresult{imagen y versión} \\

SageMath &
\todoresult{versión} \\

Python &
\todoresult{versión} \\

Biblioteca criptográfica &
\todoresult{nombre y versión} \\

Código del experimento &
\todoresult{commit Git evaluado} \\

Fecha de ejecución &
\todoresult{fecha} \\

\bottomrule
\end{tabularx}
\end{table}

Todos los candidatos fueron ejecutados sobre el mismo hardware y utilizando el mismo
mecanismo de instrumentación. Durante la toma principal de resultados se evitó ejecutar
aplicaciones ajenas al experimento que pudieran introducir carga significativa sobre
el sistema.


% ============================================================
\subsection{Configuración criptográfica}
\label{sec:crypto-configuration}

Los parámetros criptográficos empleados para cada candidato se presentan en la
Tabla~\ref{tab:crypto-parameters}.

\begin{table}[H]
\centering
\caption{Parámetros criptográficos utilizados en cada implementación.}
\label{tab:crypto-parameters}
\small
\begin{tabularx}{\textwidth}{lYYYY}
\toprule
\textbf{Protocolo} &
\textbf{Dimensión} &
\textbf{Módulo / curva} &
\textbf{Distribución / error} &
\textbf{Otros parámetros} \\
\midrule

ECDSA &
N/A &
\todoresult{curva} &
N/A &
\todoresult{hash, codificación, etc.} \\

LWE estándar &
\todoresult{$n$, $m$} &
\todoresult{$q$} &
\todoresult{$\chi_s$, $\chi_e$, $\sigma$} &
\todoresult{otros} \\

Binary-LWE &
\todoresult{$n$, $m$} &
\todoresult{$q$} &
\todoresult{secreto y error} &
\todoresult{otros} \\

Ring-LWE &
\todoresult{$n$} &
\todoresult{$q$, anillo} &
\todoresult{distribuciones} &
\todoresult{polinomio/modulus} \\

LWR &
\todoresult{$n$, $m$} &
\todoresult{$q$, $p$} &
\todoresult{distribuciones} &
\todoresult{redondeo} \\

LWE propuesto &
\todoresult{parámetros} &
\todoresult{parámetros} &
\todoresult{parámetros} &
\todoresult{parámetros} \\

\bottomrule
\end{tabularx}
\end{table}

Los parámetros fueron seleccionados buscando un nivel de seguridad comparable entre
los candidatos. Cuando no fue posible establecer una equivalencia exacta, esta
limitación se indica explícitamente en la Sección~\ref{sec:security-evaluation}.


% ============================================================
\subsection{Diseño de los experimentos}
\label{sec:experimental-design}

\subsubsection{Número de ejecuciones}

Cada operación fue ejecutada:

\begin{equation}
    N = \todoresult{número de iteraciones}
\end{equation}

veces por configuración.

Antes de iniciar la toma de datos se realizaron:

\begin{equation}
    N_{\text{warm-up}} =
    \todoresult{número de iteraciones de calentamiento}
\end{equation}

ejecuciones de calentamiento que no fueron incluidas en los resultados finales.

\subsubsection{Estadísticos reportados}

Para evitar representar el rendimiento mediante una única medición, para cada operación
se calcularon los siguientes estadísticos:

\begin{itemize}
    \item media aritmética;
    \item mediana;
    \item desviación estándar;
    \item mínimo;
    \item máximo;
    \item percentil 95;
    \item percentil 99.
\end{itemize}



% ============================================================
\subsection{Métricas evaluadas}
\label{sec:metrics}

Las métricas utilizadas se agrupan en siete dimensiones principales:

\begin{enumerate}
    \item rendimiento computacional;
    \item consumo de recursos;
    \item comunicación;
    \item seguridad;
    \item complejidad de implementación;
    \item características operacionales y experiencia de usuario;
    \item madurez y facilidad de despliegue.
\end{enumerate}


% ============================================================
\subsubsection{Rendimiento computacional}

Para cada protocolo se mide el tiempo requerido por:

\begin{itemize}
    \item generación de claves;
    \item generación del reto;
    \item generación de la respuesta;
    \item verificación;
    \item autenticación completa.
\end{itemize}

Las mediciones se reportan en milisegundos.

Adicionalmente se calcula el throughput:

\begin{equation}
    R =
    \frac{N_{\text{auth}}}
         {T_{\text{total}}},
\end{equation}

expresado en autenticaciones por segundo.


% ============================================================
\subsubsection{Consumo de memoria y almacenamiento}

Para cada implementación se estudian:

\begin{itemize}
    \item memoria RAM promedio;
    \item memoria RAM máxima;
    \item tamaño de la clave pública;
    \item tamaño de la clave privada;
    \item tamaño de parámetros públicos adicionales;
    \item estado temporal requerido durante una autenticación;
    \item almacenamiento persistente requerido por usuario.
\end{itemize}

Los tamaños de estructuras criptográficas se calculan utilizando su representación
serializada real y no únicamente su representación matemática abstracta.


% ============================================================
\subsubsection{Comunicación}

El volumen de comunicación se calcula utilizando los bytes efectivamente serializados
por la aplicación.

Para una autenticación:

\begin{equation}
    B_{\text{auth}}
    =
    B_{\text{request}}
    +
    B_{\text{challenge}}
    +
    B_{\text{response}}
    +
    B_{\text{metadata}}.
\end{equation}

Además se registra:

\begin{itemize}
    \item número de mensajes intercambiados;
    \item número de viajes de ida y vuelta (\textit{round trips});
    \item tamaño máximo de un mensaje;
    \item cantidad total de bytes transmitidos.
\end{itemize}


% ============================================================
\subsubsection{Seguridad}

La evaluación de seguridad considera:

\begin{itemize}
    \item problema matemático subyacente;
    \item parámetros utilizados;
    \item estimación de seguridad clásica;
    \item estimación de seguridad frente a adversarios cuánticos;
    \item ataques conocidos relevantes;
    \item sensibilidad a una selección incorrecta de parámetros;
    \item requisitos de aleatoriedad;
    \item probabilidad de fallo cuando sea aplicable;
    \item vulnerabilidades derivadas de la implementación.
\end{itemize}


% ============================================================
\subsubsection{Complejidad de implementación}

La dificultad de implementación no se asigna únicamente mediante una valoración
subjetiva. Se consideran los siguientes factores:

\begin{itemize}
    \item cantidad de primitivas criptográficas necesarias;
    \item cantidad de parámetros configurables;
    \item líneas de código asociadas con la implementación criptográfica;
    \item número y naturaleza de dependencias externas;
    \item necesidad de aritmética especializada;
    \item complejidad del muestreo aleatorio;
    \item complejidad de serialización;
    \item manejo de estado;
    \item disponibilidad de bibliotecas consolidadas.
\end{itemize}


% ============================================================
\subsubsection{Características operacionales}

Para estudiar las implicaciones en un sistema real se consideran:

\begin{itemize}
    \item número de pasos requeridos durante el registro;
    \item número de pasos requeridos durante una autenticación;
    \item necesidad de almacenar secretos en el dispositivo del usuario;
    \item necesidad de servidor en línea;
    \item posibilidad de verificación fuera de línea;
    \item necesidad de mantener estado por usuario;
    \item procedimiento de renovación de claves;
    \item procedimiento ante compromiso de credenciales;
    \item requerimientos de infraestructura adicionales.
\end{itemize}



\subsection{Resultados de rendimiento}
\label{sec:performance-results}

\subsubsection{Tiempo de generación de claves}

\begin{table}[H]
\centering
\caption{Tiempo de generación de claves.}
\label{tab:keygen-time}
\small
\begin{tabular}{lrrrrrr}
\toprule
\textbf{Protocolo} &
\textbf{Media} &
\textbf{Mediana} &
\textbf{Desv.} &
\textbf{P95} &
\textbf{P99} &
\textbf{Máx.} \\
&
\multicolumn{6}{c}{\textbf{ms}} \\
\midrule

ECDSA       & \todoresult{} & \todoresult{} & \todoresult{} & \todoresult{} & \todoresult{} & \todoresult{} \\
LWE         & \todoresult{} & \todoresult{} & \todoresult{} & \todoresult{} & \todoresult{} & \todoresult{} \\
Binary-LWE  & \todoresult{} & \todoresult{} & \todoresult{} & \todoresult{} & \todoresult{} & \todoresult{} \\
Ring-LWE    & \todoresult{} & \todoresult{} & \todoresult{} & \todoresult{} & \todoresult{} & \todoresult{} \\
LWR         & \todoresult{} & \todoresult{} & \todoresult{} & \todoresult{} & \todoresult{} & \todoresult{} \\
LWE propio  & \todoresult{} & \todoresult{} & \todoresult{} & \todoresult{} & \todoresult{} & \todoresult{} \\

\bottomrule
\end{tabular}
\end{table}


% ============================================================
\subsubsection{Tiempo de autenticación}

\begin{table}[H]
\centering
\caption{Tiempo promedio por etapa del proceso de autenticación.}
\label{tab:authentication-time}
\small
\begin{tabularx}{\textwidth}{lCCCCC}
\toprule
\textbf{Protocolo} &
\textbf{Challenge} &
\textbf{Response} &
\textbf{Verify} &
\textbf{Total} &
\textbf{Throughput} \\
&
\textbf{ms} &
\textbf{ms} &
\textbf{ms} &
\textbf{ms} &
\textbf{auth/s} \\
\midrule

ECDSA      & \todoresult{} & \todoresult{} & \todoresult{} & \todoresult{} & \todoresult{} \\
LWE        & \todoresult{} & \todoresult{} & \todoresult{} & \todoresult{} & \todoresult{} \\
Binary-LWE & \todoresult{} & \todoresult{} & \todoresult{} & \todoresult{} & \todoresult{} \\
Ring-LWE   & \todoresult{} & \todoresult{} & \todoresult{} & \todoresult{} & \todoresult{} \\
LWR        & \todoresult{} & \todoresult{} & \todoresult{} & \todoresult{} & \todoresult{} \\
LWE propio & \todoresult{} & \todoresult{} & \todoresult{} & \todoresult{} & \todoresult{} \\

\bottomrule
\end{tabularx}
\end{table}


% ============================================================
\subsubsection{Distribución de los tiempos}

La Figura~\ref{fig:auth-time-comparison} presenta el tiempo medio requerido para
completar una autenticación.

\begin{figure}[H]
\centering

\begin{tikzpicture}
\begin{axis}[
    ybar,
    width=0.95\textwidth,
    height=7cm,
    ylabel={Tiempo de autenticación (ms)},
    symbolic x coords={ECDSA,LWE,Binary-LWE,Ring-LWE,LWR,LWE propio},
    xtick=data,
    x tick label style={rotate=25,anchor=east},
    nodes near coords,
    ymin=0,
    enlarge x limits=0.1
]
\addplot coordinates {
    (ECDSA,0)
    (LWE,0)
    (Binary-LWE,0)
    (Ring-LWE,0)
    (LWR,0)
    (LWE propio,0)
};
\end{axis}
\end{tikzpicture}

\caption{Comparación del tiempo medio de autenticación.
Reemplazar los valores cero por los resultados experimentales.}
\label{fig:auth-time-comparison}
\end{figure}

Los resultados muestran que \todoresult{describir los resultados observados sin
interpretarlos todavía}. La diferencia entre \todoresult{protocolos relevantes} fue
de \todoresult{valor}, mientras que \todoresult{resultado adicional}.


% ============================================================
\subsection{Resultados de consumo de memoria}
\label{sec:memory-results}

\begin{table}[H]
\centering
\caption{Consumo de memoria durante las operaciones criptográficas.}
\label{tab:memory-results}
\small
\begin{tabularx}{\textwidth}{lCCCC}
\toprule
\textbf{Protocolo} &
\textbf{RAM media} &
\textbf{RAM pico} &
\textbf{Estado temporal} &
\textbf{Estado persistente} \\
&
\textbf{KB} &
\textbf{KB} &
\textbf{KB} &
\textbf{bytes/usuario} \\
\midrule

ECDSA      & \todoresult{} & \todoresult{} & \todoresult{} & \todoresult{} \\
LWE        & \todoresult{} & \todoresult{} & \todoresult{} & \todoresult{} \\
Binary-LWE & \todoresult{} & \todoresult{} & \todoresult{} & \todoresult{} \\
Ring-LWE   & \todoresult{} & \todoresult{} & \todoresult{} & \todoresult{} \\
LWR        & \todoresult{} & \todoresult{} & \todoresult{} & \todoresult{} \\
LWE propio & \todoresult{} & \todoresult{} & \todoresult{} & \todoresult{} \\

\bottomrule
\end{tabularx}
\end{table}


% ============================================================
\subsection{Tamaño de claves y estructuras criptográficas}
\label{sec:key-size-results}

\begin{table}[H]
\centering
\caption{Tamaño serializado de los principales objetos criptográficos.}
\label{tab:key-sizes}
\small
\begin{tabularx}{\textwidth}{lCCCC}
\toprule
\textbf{Protocolo} &
\textbf{Clave pública} &
\textbf{Clave privada} &
\textbf{Parámetros públicos} &
\textbf{Total registro} \\
&
\textbf{bytes} &
\textbf{bytes} &
\textbf{bytes} &
\textbf{bytes} \\
\midrule

ECDSA      & \todoresult{} & \todoresult{} & \todoresult{} & \todoresult{} \\
LWE        & \todoresult{} & \todoresult{} & \todoresult{} & \todoresult{} \\
Binary-LWE & \todoresult{} & \todoresult{} & \todoresult{} & \todoresult{} \\
Ring-LWE   & \todoresult{} & \todoresult{} & \todoresult{} & \todoresult{} \\
LWR        & \todoresult{} & \todoresult{} & \todoresult{} & \todoresult{} \\
LWE propio & \todoresult{} & \todoresult{} & \todoresult{} & \todoresult{} \\

\bottomrule
\end{tabularx}
\end{table}


% ============================================================
\subsection{Resultados de comunicación}
\label{sec:communication-results}

\begin{table}[H]
\centering
\caption{Datos transmitidos durante una autenticación.}
\label{tab:communication}
\small
\begin{tabularx}{\textwidth}{lCCCCC}
\toprule
\textbf{Protocolo} &
\textbf{Challenge} &
\textbf{Response} &
\textbf{Metadata} &
\textbf{Total} &
\textbf{Mensajes} \\
&
\textbf{bytes} &
\textbf{bytes} &
\textbf{bytes} &
\textbf{bytes} &
\\
\midrule

ECDSA      & \todoresult{} & \todoresult{} & \todoresult{} & \todoresult{} & \todoresult{} \\
LWE        & \todoresult{} & \todoresult{} & \todoresult{} & \todoresult{} & \todoresult{} \\
Binary-LWE & \todoresult{} & \todoresult{} & \todoresult{} & \todoresult{} & \todoresult{} \\
Ring-LWE   & \todoresult{} & \todoresult{} & \todoresult{} & \todoresult{} & \todoresult{} \\
LWR        & \todoresult{} & \todoresult{} & \todoresult{} & \todoresult{} & \todoresult{} \\
LWE propio & \todoresult{} & \todoresult{} & \todoresult{} & \todoresult{} & \todoresult{} \\

\bottomrule
\end{tabularx}
\end{table}


% ============================================================
\subsection{Escalabilidad de almacenamiento}
\label{sec:storage-scalability}

A partir del almacenamiento requerido por usuario se estima el espacio necesario para
diferentes tamaños del sistema.

Si $S_u$ representa el almacenamiento persistente requerido por usuario, entonces:

\begin{equation}
    S(N) = N S_u.
\end{equation}

\begin{table}[H]
\centering
\caption{Estimación de almacenamiento según el número de usuarios.}
\label{tab:storage-scalability}
\small
\begin{tabularx}{\textwidth}{lCCCC}
\toprule
\textbf{Protocolo} &
\textbf{1 usuario} &
\textbf{1.000 usuarios} &
\textbf{100.000 usuarios} &
\textbf{1.000.000 usuarios} \\
\midrule

ECDSA      & \todoresult{} & \todoresult{} & \todoresult{} & \todoresult{} \\
LWE        & \todoresult{} & \todoresult{} & \todoresult{} & \todoresult{} \\
Binary-LWE & \todoresult{} & \todoresult{} & \todoresult{} & \todoresult{} \\
Ring-LWE   & \todoresult{} & \todoresult{} & \todoresult{} & \todoresult{} \\
LWR        & \todoresult{} & \todoresult{} & \todoresult{} & \todoresult{} \\
LWE propio & \todoresult{} & \todoresult{} & \todoresult{} & \todoresult{} \\

\bottomrule
\end{tabularx}
\end{table}


% ============================================================
\subsection{Evaluación de seguridad}
\label{sec:security-evaluation}

La seguridad de las alternativas basadas en retículos no se determina únicamente por
el nombre de la familia criptográfica. Depende de los parámetros concretos
seleccionados para cada implementación.

Por este motivo, para cada candidato se documentan los parámetros y la estimación de
seguridad correspondiente.

\begin{table}[H]
\centering
\caption{Nivel de seguridad estimado para las configuraciones evaluadas.}
\label{tab:security-levels}
\small
\begin{tabularx}{\textwidth}{lYCCC}
\toprule
\textbf{Protocolo} &
\textbf{Problema base} &
\textbf{Seguridad clásica} &
\textbf{Seguridad cuántica} &
\textbf{Método de estimación} \\
\midrule

ECDSA &
ECDLP &
\todoresult{} &
\todoresult{} &
\todoresult{} \\

LWE &
LWE &
\todoresult{} &
\todoresult{} &
\todoresult{} \\

Binary-LWE &
LWE con secreto binario &
\todoresult{} &
\todoresult{} &
\todoresult{} \\

Ring-LWE &
Ring-LWE &
\todoresult{} &
\todoresult{} &
\todoresult{} \\

LWR &
LWR &
\todoresult{} &
\todoresult{} &
\todoresult{} \\

LWE propio &
\todoresult{} &
\todoresult{} &
\todoresult{} &
\todoresult{} \\

\bottomrule
\end{tabularx}
\end{table}

Cuando se utilicen estimadores de seguridad, deberán registrarse también la versión de
la herramienta, sus parámetros de entrada y los modelos de costo utilizados.


% ============================================================

\subsection{Análisis de ataques}
\label{sec:attack-analysis}

La Tabla~\ref{tab:attack-analysis} resume las principales familias de ataques relevantes
para cada protocolo. La presencia de una familia de ataque no implica necesariamente
que el protocolo sea vulnerable bajo los parámetros empleados; se documenta para
identificar las condiciones que deben satisfacerse para conservar el nivel de
seguridad esperado.

\begin{longtable}{p{2.5cm}p{3.0cm}p{3.6cm}p{4.1cm}}
\caption{Ataques relevantes y mecanismos de mitigación.}
\label{tab:attack-analysis}
\\

\toprule
\textbf{Protocolo} &
\textbf{Ataque / riesgo} &
\textbf{Condición relevante} &
\textbf{Mitigación / consideración} \\
\midrule
\endfirsthead

\toprule
\textbf{Protocolo} &
\textbf{Ataque / riesgo} &
\textbf{Condición relevante} &
\textbf{Mitigación / consideración} \\
\midrule
\endhead

ECDSA &
\todoresult{ataque} &
\todoresult{condición} &
\todoresult{mitigación} \\

ECDSA &
\todoresult{ataque} &
\todoresult{condición} &
\todoresult{mitigación} \\

LWE &
\todoresult{ataque} &
\todoresult{condición} &
\todoresult{mitigación} \\

LWE &
\todoresult{ataque} &
\todoresult{condición} &
\todoresult{mitigación} \\

Binary-LWE &
\todoresult{ataque} &
\todoresult{condición} &
\todoresult{mitigación} \\

Ring-LWE &
\todoresult{ataque} &
\todoresult{condición} &
\todoresult{mitigación} \\

LWR &
\todoresult{ataque} &
\todoresult{condición} &
\todoresult{mitigación} \\

LWE propio &
\todoresult{ataque} &
\todoresult{condición} &
\todoresult{mitigación} \\

\bottomrule
\end{longtable}


% ============================================================
\subsection{Seguridad de implementación}
\label{sec:implementation-security}

Además de la seguridad matemática de cada construcción, se analizaron riesgos
relacionados con su implementación.

\begin{table}[H]
\centering
\caption{Aspectos de seguridad operacional de las implementaciones.}
\label{tab:implementation-security}
\small
\begin{tabularx}{\textwidth}{lCCCCY}
\toprule
\textbf{Protocolo} &
\textbf{RNG} &
\textbf{Nonce} &
\textbf{Side channels} &
\textbf{Parámetros} &
\textbf{Observaciones} \\
\midrule

ECDSA      & \todoresult{} & \todoresult{} & \todoresult{} & \todoresult{} & \todoresult{} \\
LWE        & \todoresult{} & \todoresult{} & \todoresult{} & \todoresult{} & \todoresult{} \\
Binary-LWE & \todoresult{} & \todoresult{} & \todoresult{} & \todoresult{} & \todoresult{} \\
Ring-LWE   & \todoresult{} & \todoresult{} & \todoresult{} & \todoresult{} & \todoresult{} \\
LWR        & \todoresult{} & \todoresult{} & \todoresult{} & \todoresult{} & \todoresult{} \\
LWE propio & \todoresult{} & \todoresult{} & \todoresult{} & \todoresult{} & \todoresult{} \\

\bottomrule
\end{tabularx}
\end{table}


% ============================================================

\subsection{Requisitos operacionales y experiencia de usuario}
\label{sec:operational-results}

\begin{table}[H]
\centering
\caption{Comparación de requisitos operacionales.}
\label{tab:operational-comparison}
\small
\begin{tabularx}{\textwidth}{lCCCCCY}
\toprule
\textbf{Protocolo} &
\textbf{Registro} &
\textbf{Estado} &
\textbf{Servidor online} &
\textbf{Secreto cliente} &
\textbf{Verif. offline} &
\textbf{Observaciones} \\
\midrule

ECDSA      & \todoresult{} & \todoresult{} & \todoresult{} & \todoresult{} & \todoresult{} & \todoresult{} \\
LWE        & \todoresult{} & \todoresult{} & \todoresult{} & \todoresult{} & \todoresult{} & \todoresult{} \\
Binary-LWE & \todoresult{} & \todoresult{} & \todoresult{} & \todoresult{} & \todoresult{} & \todoresult{} \\
Ring-LWE   & \todoresult{} & \todoresult{} & \todoresult{} & \todoresult{} & \todoresult{} & \todoresult{} \\
LWR        & \todoresult{} & \todoresult{} & \todoresult{} & \todoresult{} & \todoresult{} & \todoresult{} \\
LWE propio & \todoresult{} & \todoresult{} & \todoresult{} & \todoresult{} & \todoresult{} & \todoresult{} \\

\bottomrule
\end{tabularx}
\end{table}

Desde la perspectiva del usuario final, \todoresult{describir cuáles diferencias son
visibles para el usuario y cuáles quedan ocultas por la aplicación}.


% ============================================================
\subsection{Dependencias de confianza y arquitectura}
\label{sec:trust-model}

En lugar de tratar la centralización como una propiedad binaria, se analizan
explícitamente las dependencias de confianza de cada construcción.

\begin{table}[H]
\centering
\caption{Dependencias de confianza e infraestructura.}
\label{tab:trust-model}
\small
\begin{tabularx}{\textwidth}{lCCCCYY}
\toprule
\textbf{Protocolo} &
\textbf{Autoridad} &
\textbf{PKI} &
\textbf{Estado servidor} &
\textbf{Secreto servidor} &
\textbf{Compromiso servidor} &
\textbf{Observaciones} \\
\midrule

ECDSA      & \todoresult{} & \todoresult{} & \todoresult{} & \todoresult{} & \todoresult{} & \todoresult{} \\
LWE        & \todoresult{} & \todoresult{} & \todoresult{} & \todoresult{} & \todoresult{} & \todoresult{} \\
Binary-LWE & \todoresult{} & \todoresult{} & \todoresult{} & \todoresult{} & \todoresult{} & \todoresult{} \\
Ring-LWE   & \todoresult{} & \todoresult{} & \todoresult{} & \todoresult{} & \todoresult{} & \todoresult{} \\
LWR        & \todoresult{} & \todoresult{} & \todoresult{} & \todoresult{} & \todoresult{} & \todoresult{} \\
LWE propio & \todoresult{} & \todoresult{} & \todoresult{} & \todoresult{} & \todoresult{} & \todoresult{} \\

\bottomrule
\end{tabularx}
\end{table}

% ============================================================
\subsection{Comparación global}
\label{sec:global-comparison}

La Tabla~\ref{tab:global-summary} resume las principales métricas cuantitativas.
Esta tabla no pretende producir un ranking absoluto entre protocolos, sino permitir
identificar visualmente los compromisos existentes entre las diferentes alternativas.

\begin{table}[H]
\centering
\caption{Resumen cuantitativo de los resultados principales.}
\label{tab:global-summary}
\resizebox{\textwidth}{!}{
\begin{tabular}{lrrrrrr}
\toprule

\textbf{Protocolo} &
\textbf{Auth. (ms)} &
\textbf{RAM pico} &
\textbf{Comunicación} &
\textbf{Clave pública} &
\textbf{Clave privada} &
\textbf{Seguridad estimada} \\

&
&
\textbf{KB} &
\textbf{bytes} &
\textbf{bytes} &
\textbf{bytes} &
\textbf{bits} \\

\midrule

ECDSA
& \todoresult{}
& \todoresult{}
& \todoresult{}
& \todoresult{}
& \todoresult{}
& \todoresult{} \\

LWE
& \todoresult{}
& \todoresult{}
& \todoresult{}
& \todoresult{}
& \todoresult{}
& \todoresult{} \\

Binary-LWE
& \todoresult{}
& \todoresult{}
& \todoresult{}
& \todoresult{}
& \todoresult{}
& \todoresult{} \\

Ring-LWE
& \todoresult{}
& \todoresult{}
& \todoresult{}
& \todoresult{}
& \todoresult{}
& \todoresult{} \\

LWR
& \todoresult{}
& \todoresult{}
& \todoresult{}
& \todoresult{}
& \todoresult{}
& \todoresult{} \\

LWE propio
& \todoresult{}
& \todoresult{}
& \todoresult{}
& \todoresult{}
& \todoresult{}
& \todoresult{} \\

\bottomrule
\end{tabular}
}
\end{table}


% ============================================================
\subsection{Análisis multidimensional}
\label{sec:multidimensional-analysis}

Una única métrica no permite determinar la conveniencia de una construcción
criptográfica. Un protocolo puede reducir el tiempo de ejecución a costa de incrementar
el tamaño de las claves, mientras que otro puede reducir la comunicación aumentando
la complejidad computacional.

Por esta razón, los resultados se analizan como un conjunto de compromisos entre:

\begin{equation}
    (\text{seguridad},
     \text{tiempo},
     \text{memoria},
     \text{comunicación},
     \text{complejidad}).
\end{equation}

\subsubsection{Tiempo frente a comunicación}

\begin{figure}[H]
\centering
\fbox{
\parbox[c][6cm][c]{0.90\textwidth}{
\centering
\todoresult{Insertar gráfica de dispersión: \\
eje X = tiempo de autenticación; \\
eje Y = bytes transmitidos; \\
tamaño del punto = RAM pico.}
}}
\caption{Relación entre tiempo de autenticación, volumen de comunicación y memoria.}
\label{fig:time-bandwidth}
\end{figure}


% ============================================================
\subsubsection{Seguridad frente a rendimiento}

\begin{figure}[H]
\centering
\fbox{
\parbox[c][6cm][c]{0.90\textwidth}{
\centering
\todoresult{Insertar gráfica: \\
eje X = seguridad estimada; \\
eje Y = tiempo de autenticación.}
}}
\caption{Relación entre seguridad estimada y costo computacional.}
\label{fig:security-performance}
\end{figure}


\subsection{Discusión}
\label{sec:benchmark-discussion}

Los experimentos muestran que no existe necesariamente una alternativa que optimice
simultáneamente todas las dimensiones estudiadas.

Desde el punto de vista de rendimiento,
\todoresult{hallazgo}.

Respecto al consumo de memoria,
\todoresult{hallazgo}.

En términos de comunicación,
\todoresult{hallazgo}.

Respecto a la seguridad,
\todoresult{hallazgo}.

En términos de complejidad de implementación y despliegue,
\todoresult{hallazgo}.

Estos resultados evidencian un compromiso entre
\todoresult{trade-offs principales}.


% ============================================================
\subsubsection{Comparación con ECDSA}
\label{subsec:comparison-ecdsa}

ECDSA constituye el punto de referencia utilizado durante la evaluación.

Frente a esta alternativa, las construcciones basadas en retículos presentan
\todoresult{ventajas observadas}, mientras que presentan costos adicionales asociados
con \todoresult{desventajas observadas}.

En particular,
\todoresult{resultado cuantitativo importante}.


% ============================================================
\subsubsection{Comparación entre las variantes basadas en LWE}
\label{subsec:lwe-comparison}

Las variantes evaluadas muestran diferencias significativas en
\todoresult{métricas relevantes}.

LWE estándar presentó
\todoresult{resultado}.

Binary-LWE presentó
\todoresult{resultado}.

Ring-LWE presentó
\todoresult{resultado}.

LWR presentó
\todoresult{resultado}.

Estas diferencias permiten observar cómo modificaciones en la representación del
secreto, la estructura algebraica o el mecanismo utilizado para introducir
incertidumbre afectan los costos de una implementación real.


% ============================================================
\subsubsection{Evaluación del protocolo propuesto}
\label{subsec:proposed-evaluation}

El protocolo propuesto debe analizarse utilizando exactamente los mismos criterios que
las demás alternativas.

Los resultados muestran que
\todoresult{hallazgo principal}.

Su mayor ventaja observada corresponde a
\todoresult{ventaja}.

Su principal costo o limitación corresponde a
\todoresult{limitación}.

Frente a las demás variantes LWE,
\todoresult{comparación}.

Frente a ECDSA,
\todoresult{comparación}.

Por tanto, los resultados sugieren que el protocolo
\todoresult{interpretación prudente basada exclusivamente en los resultados obtenidos}.


% ============================================================
\subsection{Amenazas a la validez}
\label{sec:threats-validity}

Los resultados presentados deben interpretarse considerando las siguientes amenazas a
la validez.

\subsubsection{Validez interna}

Los tiempos de ejecución pueden verse afectados por procesos del sistema operativo,
planificación del procesador, cachés y otros factores externos. Para reducir este
efecto se utilizaron múltiples repeticiones, iteraciones de calentamiento y un entorno
experimental controlado.

\todoresult{Agregar otras amenazas internas identificadas.}


\subsubsection{Validez externa}

Los resultados corresponden al hardware, sistema operativo, versiones de software y
parámetros descritos en la Sección~\ref{sec:experimental-environment}. Por esta razón,
los valores absolutos pueden variar al utilizar hardware diferente.

Sin embargo, el diseño experimental busca permitir la reproducción del procedimiento
en otros entornos.

\todoresult{Agregar otras amenazas externas.}


\subsubsection{Validez de constructo}

Las métricas utilizadas representan diferentes dimensiones de eficiencia y
desplegabilidad, pero ninguna métrica individual captura completamente la calidad de
un protocolo criptográfico.

En particular, indicadores como líneas de código, cantidad de dependencias o número de
parámetros deben interpretarse únicamente como aproximaciones a la complejidad de
implementación y no como medidas absolutas de dificultad.


\subsubsection{Limitaciones de la estimación de seguridad}

Las estimaciones de seguridad dependen de los modelos de ataque, parámetros y
herramientas disponibles al momento de realizar el estudio.

\todoresult{Describir herramienta de estimación utilizada, versión, supuestos y
limitaciones.}


% ============================================================
\subsection{Reproducibilidad}
\label{sec:reproducibility}

Con el propósito de facilitar la reproducción de los experimentos, el repositorio del
proyecto conserva:

\begin{itemize}
    \item las implementaciones evaluadas;
    \item los parámetros criptográficos utilizados;
    \item los scripts de benchmarking;
    \item la configuración del entorno de ejecución;
    \item las versiones de las dependencias;
    \item las migraciones o estructura de almacenamiento;
    \item los resultados crudos de cada ejecución;
    \item los procedimientos utilizados para generar las métricas agregadas.
\end{itemize}

El experimento correspondiente a este documento fue ejecutado utilizando el commit:

\begin{center}
    \texttt{\todoresult{hash del commit}}
\end{center}

y la configuración experimental:

\begin{center}
    \texttt{\todoresult{identificador de configuración}}
\end{center}



\section{Conclusiones}

A partir de los experimentos realizados se identificaron los siguientes hallazgos
principales:

\begin{enumerate}
    \item \todoresult{hallazgo principal sobre rendimiento}.

    \item \todoresult{hallazgo principal sobre memoria}.

    \item \todoresult{hallazgo principal sobre comunicación}.

    \item \todoresult{hallazgo principal sobre seguridad}.

    \item \todoresult{hallazgo principal sobre implementación y despliegue}.

    \item \todoresult{hallazgo principal sobre el protocolo propuesto}.
\end{enumerate}

En conjunto, los resultados permiten caracterizar las diferentes alternativas no
mediante una clasificación única, sino a partir de los compromisos que cada una
establece entre seguridad, costo computacional, memoria, comunicación y complejidad
operacional.

Esta caracterización proporciona la base experimental para determinar en qué contextos
resulta apropiado utilizar cada familia de protocolos y permite evaluar de forma
homogénea la propuesta desarrollada en este trabajo.

% 5. Bibliografía en Estilo IEEE
\newpage
\bibliographystyle{IEEEtran}
\bibliography{referencias.bib}

\end{document}